\documentclass[10pt]{article}
\usepackage[margin=0.8in]{geometry}
\usepackage[T1]{fontenc}
\usepackage{amsmath,amssymb,amsthm}
\usepackage{multicol}
\usepackage{hyperref}
\newtheorem{theorem}{Theorem}
\newtheorem{lemma}{Lemma}
\newtheorem{corollary}{Corollary}
\title{Cube-Sum Permutations: Unimodular Block Algebra and 57 Finite Shapes}
\author{Qichao Wang\\{\small Hebei University of Technology, Tianjin, China}\\[0.4em]Daoyu Dong\\{\small University of Electronic Science and Technology of China, Chengdu, China}}
\date{October 2026}

\begin{document}
\maketitle

\begin{abstract}
Let $P(n)$ count permutations $p$ of $[1,n]$
such that every $i+p(i)$ is a perfect cube.
We prove $P(n)>0$ for every $n\ge173$
from six explicit block-rearrangement families,
$57$ finite cube-word/output-order shapes, and
elementary cubic suffix reflection.
A general algebraic theorem proves that the
block-order matrix has determinant $0$ or $1$.
For each of the $57$ shapes the determinant is $1$,
and a single Cramer-determinant formula yields
strictly positive integral block lengths.
The permutation then follows symbolically,
without separately verifying $57$ endpoint
multiset equalities. Every such shape also admits
a uniform infinite cubic dilation
$N(h)=(N+\varepsilon)h^3-\varepsilon$.
A symbolic cubic-gap induction completes
the existence tail for all $n\ge365$.
The $57$ shape choices and positivity checks
remain finite casework; this is an algebraic
parameterization, not a case-free existence proof.
\end{abstract}

\section{Two elementary lemmas}

\begin{lemma}[Boundary cancellation]\label{lem:balance}
Let $0=b_0<b_1<\cdots<b_t=n$ and let
$c_1,\ldots,c_t$ be positive integer cubes.
Set $L_j=c_j-b_j$ and $R_j=c_j-b_{j-1}$.
Assume $1\le L_j<R_j\le n+1$.
Then $p(i)=c_j-i$ for $b_{j-1}<i\le b_j$
is a cube-sum permutation if and only if
\[
\boxed{\{1,R_1,\ldots,R_t\}_{\mathrm{multi}}
 =\{n+1,L_1,\ldots,L_t\}_{\mathrm{multi}}.}
\]
\end{lemma}
\begin{proof}
Each input block maps bijectively onto the integer
interval $[L_j,R_j)$. At an integer $x$, the jump
in the total interval-coverage count is the
multiplicity of $x$ as a left endpoint minus
its multiplicity as a right endpoint.
The displayed identity makes that jump $1$
at $x=1$, $-1$ at $x=n+1$, and zero everywhere
else. Thus coverage is exactly one at
every $x\in[1,n]$. Reversing this argument
proves necessity.
\end{proof}

\begin{lemma}[Growing cubic suffix]\label{lem:suffix}
Suppose $0\le m<n$, $m+n+1=c$ is a cube,
and a cube-sum permutation on $[1,m]$ exists
(the $m=0$ permutation is empty).
Keep the existing permutation on $[1,m]$
and set $p(i)=c-i$ for $m<i\le n$.
Then $p$ is a cube-sum permutation on $[1,n]$.
\end{lemma}
\begin{proof}
The new reflection is a bijection from
$[m+1,n]$ to itself, and each new sum is $c$.
\end{proof}

\section{Six algebraic block families and 57 exceptional identities}

\begin{lemma}[Reflected-block rearrangement]\label{lem:rearrange}
Choose positive input-block lengths
$\ell_1,\ldots,\ell_t$ of sum $n$ and an arbitrary
output order $\pi$ of these $t$ blocks.
Let $A_j$ be the sum of input lengths before
block $j$, and let $B_j$ be the sum of output
lengths before block $j$ in order $\pi$.
If all numbers
\[
  C_j=1+A_j+B_j+\ell_j
\]
are perfect cubes, then on input block $j$
the reflection $p(i)=C_j-i$ gives a cube-sum
permutation of $[1,n]$.
\end{lemma}
\begin{proof}
The input block is $[A_j+1,A_j+\ell_j]$,
which maps exactly to the prescribed output
block $[B_j+1,B_j+\ell_j]$.
Those output blocks partition $[1,n]$
by construction. Each index sum is $C_j$.
\end{proof}

\paragraph{Family I (four blocks).}
Let $U<V$ be integer cubes and put
\[
a=2U-V-1,\qquad b=2V-3U+1.
\]
If $a,b>0$, use input lengths $(a,b,a,b)$
and output order $(3,4,1,2)$.
Lemma~\ref{lem:rearrange} gives the cubic
sums $(U,V,U,V)$ and length $2(V-U)$.
The pairs $(U,V)=(64,125),(125,216),(216,343)$
yield lengths $122,182,254$ respectively.

\paragraph{Family II (five blocks).}
For cubes $X<Y<Z$, define
\[
a=2Y-Z-1,\quad
b=X+2Z-4Y+1,\quad c=Y-X.
\]
If $a,b,c>0$, the lengths $(a,b,c,a,c)$
in output order $(4,2,5,1,3)$ produce
sums $(Y,X,Z,Y,Z)$ and total length
$2Y-X-1$. With $(Y,Z)=(64,125)$ and
$X=27,8$, the results are $n=100,119$.

\paragraph{Family III (six blocks).}
For integer cubes $X,W,Y,Z$, put
\[
\begin{aligned}
a&=2Y-Z-1,& b&=X-2a-1,\\
c&=W+2a-2X+1,& d&=Y+X-W-2a-1.
\end{aligned}
\]
If these are positive, choose lengths
$(a,b,c,d,a,d)$ and output order
$(5,2,3,6,1,4)$. The cubic sums are
$(Y,X,W,Z,Y,Z)$ and the length is
$2Z-2Y+X-W$. Taking $(Y,Z)=(64,125)$,
the choices $(X,W)=(8,64),(27,64),(8,27)$
give $n=66,85,103$, respectively.

\paragraph{Family IV (twelve blocks and three binary switches).}
Let $X,T,Q$ be integers and put
\[
\begin{aligned}
u&=X-7,&v&=T-2X+7,&w&=X+27-T,\\
z&=Q+2T-2X-121,&h&=X+152-T-Q.
\end{aligned}
\]
Assume $u,v,z\ge0$ and $w,h>0$.
Use input lengths
\[
(3,u,v,w,3,w,z,h,3,27,3,h)
\]
and output block order
\[
(11,2,3,6,9,4,7,12,5,10,1,8).
\]
The reflected-block lemma computes the twelve
nominal index-image sums as
\[
(125,X,T,64,125,64,Q,216,125,216,125,216).
\]
When the second, third, or seventh block is nonempty,
its nominal sum $X$, $T$, or $Q$, respectively,
must be a cube. Any zero-length block is omitted
from both orders and imposes no cube condition.
The total length is
\[
n=276+X-T-Q.
\]
Take $X\in\{7,8\}$, choose
$T\in\{2X-7,27\}$, and independently choose
$Q\in\{121-2T+2X,125\}$. All eight resulting
length choices satisfy the positivity and cube
conditions. The $2^3$ distinct lengths are
\[
\boxed{131,132,150,151,155,156,174,175}.
\]
This is one symbolic rearrangement construction
with three optional empty blocks, not eight
separate finite endpoint certificates.

\paragraph{Family V (15 blocks, two constrained switches).}
Take $\alpha,\beta\in\{0,1\}$ with
$\alpha\le\beta$ and
$t\in\{0,7+4\beta+2\alpha\}$.
Define the ordered input block lengths
\[
\begin{gathered}
\ell_1=2\beta,\quad
\ell_2=7-\alpha-4\beta-\alpha\beta,\quad
\ell_3=\alpha,\quad\ell_4=\alpha\beta,\\
\ell_5=13+4\beta,\quad\ell_6=\alpha,\quad
\ell_7=23-4\beta-\alpha,\quad\ell_8=13+4\beta,\\
\ell_9=48-\alpha-45\beta+\alpha\beta,\quad
\ell_{10}=2\beta,\quad\ell_{11}=\beta(39-\alpha),\\
\ell_{12}=t,\quad\ell_{13}=30+\alpha-t,\quad
\ell_{14}=13+4\beta,\quad
\ell_{15}=91+\alpha-t.
\end{gathered}
\]
The output block order is
\[
(10,4,6,2,8,3,11,1,9,14,7,12,15,5,13).
\]
All lengths are nonnegative, with empty blocks omitted.
Because $\alpha\beta=\alpha$, the nominal reflected
sums are
\[
\begin{gathered}
21+43\beta,\ 8,\ 28-\alpha,\ 8,\ 216,\\
21+2\alpha+4\beta,\ 125,\ 64,\ 125,\ 105-41\beta,\\
125,\ 209-2\alpha-4\beta+t,\ 343,\ 216,\ 343.
\end{gathered}
\]
Every \emph{nonempty} block has a cubic sum:
blocks $1,10,11$ are used only if $\beta=1$;
blocks $3,4,6$ are used only if
$\alpha=\beta=1$; and block $12$ is empty
if $t=0$, otherwise its sum is $216$.
The total length is
\[
\boxed{n=238+\alpha+2\beta-t},
\]
so the three allowed pairs
$(\alpha,\beta)=(0,0),(0,1),(1,1)$
with their two possible $t$ values give
\[
\boxed{228,229,231,238,240,241}.
\]
This is a symbolic cubic-sum construction;
it does not require verifying six separate endpoint lists.

\paragraph{Family VI (ten blocks, nested binary degeneration).}
Let $\alpha,\beta\in\{0,1\}$ with $\beta\le\alpha$.
Use input block lengths
\[
\begin{gathered}
(7,\ 12\alpha,\ 159-24\alpha-99\beta,\
14\beta,\ 28\beta,\ 63\beta,\\
10+24\alpha-6\beta,\ 8\beta,\
159-24\alpha-107\beta,\
10+24\alpha+93\beta)
\end{gathered}
\]
and output order $(1,2,8,5,6,9,4,10,3,7)$.
Direct substitution into the reflected-block lemma
gives the ten nominal sums
\[
\begin{gathered}
8,\ 8+19\alpha,\ 343,\
167-12\alpha+61\beta,\
167-12\alpha-30\beta,\\
167-12\alpha+61\beta,\ 512,\
177+12\alpha+27\beta,\ 343,\ 512.
\end{gathered}
\]
Block $2$ is nonempty only when $\alpha=1$
and has cubic sum $27$; blocks $4,5,6,8$
are nonempty only when $\beta=\alpha=1$
and have sums $216,125,216,216$. Other
nonempty blocks have constant cube sums;
all empty blocks are omitted. Since
\[
n=345+12\alpha-6\beta,
\]
the three admissible choices give the
lengths $345,357,351$, using one
parametric construction rather than
three independent endpoint lists.

\subsection{Universal unimodular determinant construction}
\label{subsec:unimodular-factory}
The remaining $57$ integer witnesses admit a common
algebraic formula, without separately testing
their endpoint cancellation identities.
For an output-block order $\pi$ on $t$ labels,
write $\rho(i)=\pi^{-1}(i)$ and set
\[
M_{ji}=[i<j]+[\rho(i)<\rho(j)]+[i=j],
\quad 1\le i,j\le t.
\]
Here brackets denote indicator functions.

\begin{lemma}[Unimodular block-order matrices]
\label{lem:unimodular-block}
For every finite output-block order $\pi$,
$\det M\in\{0,1\}$.
\end{lemma}
\begin{proof}
With $\boldsymbol1$ denoting the all-ones
vector, write $M=\boldsymbol1\boldsymbol1^T+K$.
The matrix $K$ is skew-symmetric, and for
$i<j$ its $(i,j)$ entry is $-1$ exactly
when $\rho(i)<\rho(j)$, and $0$ otherwise.
It is therefore an oriented chord-intersection
matrix: place first chord endpoints in
input order around one arc of a circle,
and second endpoints in output order around
the complementary arc.

Recall that an even-order chord-intersection
matrix has Pfaffian $0$ or $\pm1$.
Indeed, thickening the chord diagram into an
oriented ribbon surface, the chord loops form an
integral homology basis. Their intersection form
is integrally congruent to a direct sum of
unimodular symplectic pairs
$\left(\begin{smallmatrix}0&1\\-1&0\end{smallmatrix}\right)$
and zero blocks from the boundary radical.
Its full Pfaffian is thus $0$ or $\pm1$.

If $t$ is even, $\operatorname{adj}(K)$ is
skew-symmetric; the rank-one determinant
identity gives
\[
\det M=\det K+
\boldsymbol1^T\operatorname{adj}(K)\boldsymbol1
=\operatorname{Pf}(K)^2\in\{0,1\}.
\]
If $t$ is odd, $\det K=0$. Adjoin a new
chord crossing all original chords, whose
skew matrix is
\[
H=\begin{pmatrix}
0&\boldsymbol1^T\\
-\boldsymbol1&K
\end{pmatrix}.
\]
The same rank-one and bordered determinant
identities give
\[
\det M=\boldsymbol1^T\operatorname{adj}(K)
\boldsymbol1=\det H=\operatorname{Pf}(H)^2
\in\{0,1\}.
\]
\end{proof}

\begin{theorem}[Symbolic determinant families]
\label{thm:unimodular-block-factory}
Fix $\pi$ with $\det M=1$.
For arbitrary integer cube roots
$k_1,\ldots,k_t\ge2$, let
$d=(k_j^3-1)_{j=1}^t$. Define
\[
\ell_j=\det M^{(j)}(d),
\]
where $M^{(j)}(d)$ is obtained from $M$
by replacing column $j$ with $d$.
If every $\ell_j>0$, then with
$N=\sum_j\ell_j$ there is a permutation
of $[1,N]$ with cube sums $k_j^3$ on
successive input blocks of these lengths.
\end{theorem}
\begin{proof}
By the preceding lemma and Cramer's rule,
$M\ell=d$ holds identically over the integers.
If $A_j,B_j$ are the input and output block
prefix lengths, this is precisely
\[
k_j^3=1+A_j+B_j+\ell_j.
\]
Reflection $p(i)=k_j^3-i$ maps the
$j$th input block bijectively to the
$j$th block in output order $\pi$.
The output blocks partition $[1,N]$.
\end{proof}

The $57$ cubic-root words and block orders
in Appendix~\ref{app:algebraic-shapes}
are therefore a collection of \emph{symbolic}
families, rather than $57$ independently
postulated endpoint identities.
For each printed root word,
integer determinant evaluation gives
$\det M=1$, $\ell_j\ge1$, and the
specified length $N$ (the smallest
$\ell_j$ across the full table equals $1$).
The separate exact
\texttt{verify\_unimodular\_factory.py}
checks these finite positivity conditions.
The $57$ parameter choices remain finite data:
this is not a case-free proof of their existence.

\begin{theorem}[Uniform cubic dilation]
\label{thm:unimodular-dilation}
Under Theorem~\ref{thm:unimodular-block-factory},
put $v=M^{-1}\boldsymbol1$ and
$\varepsilon=t\bmod2$. For every integer $h\ge1$,
\[
\ell_j(h)=\ell_j+
(h^3-1)(\ell_j+v_j)>0
\]
defines another cube-sum permutation, with
root word $(hk_1,\ldots,hk_t)$, of length
\[
\boxed{N(h)=(N+\varepsilon)h^3-\varepsilon.}
\]
\end{theorem}
\begin{proof}
Pfaffian minor expansions in the preceding
chord-diagram proof imply
$v_j\in\{-1,0,1\}$: the Cramer minors
are signed Pfaffians of chord diagrams
with a marked chord deleted, divided
by a nonzero Pfaffian of absolute value $1$.
Moreover $\sum_j v_j=\varepsilon$.
Indeed, $Mv=\boldsymbol1$ and skew symmetry
give $S^2=S$ for $S=\boldsymbol1^Tv$.
If $t$ is even, $K$ is invertible and
$v=K^{-1}\boldsymbol1$ yields $S=0$.
If $t$ is odd, invertibility of $M$ and
singularity of $K$ exclude $S=0$:
a kernel vector of $K$ has nonzero pairing
with $\boldsymbol1$. Thus $S=1$.

Since $\ell_j\ge1$ and $v_j\ge-1$,
all displayed scaled lengths are positive.
Finally
\[
M\ell(h)=h^3M\ell+(h^3-1)\boldsymbol1
=h^3(k_j^3)_{j=1}^t-\boldsymbol1.
\]
The determinant construction applies again.
Summing the lengths gives the formula
for $N(h)$.
\end{proof}

\begin{lemma}[Affine-word rigidity]\label{lem:rigid}
Fix a finite list of integer cubes $C_1,\ldots,C_t$.
Suppose all breakpoints $b_j(n)$ are affine
functions of $n$ with $b_0(n)=0$ and $b_t(n)=n$.
The resulting fixed-cube reflection word can
be a valid permutation for at most two distinct
values of $n$.
\end{lemma}
\begin{proof}
Every legal permutation satisfies the sum identity
\[
 \sum_{j=1}^{t} C_j\bigl(b_j(n)-b_{j-1}(n)\bigr)
 =\sum_{i=1}^n(i+p(i))=n(n+1).
\]
Its left side is affine in $n$; its right
side is a monic quadratic. The equality
cannot hold for three distinct values.
\end{proof}
The obstruction concerns \emph{fixed} cube sums
and affine endpoints only, not varying cube
sums or composition of reflection templates.

In total these \emph{six symbolic identities}
give twenty-five base lengths
\[
66,85,100,103,119,122,131,132,150,151,155,156,174,175,182,228,229,231,238,240,241,254,345,351,357
\]
without twenty-five separate endpoint-table checks.

\paragraph{Redundant 57-row endpoint atlas.}

The appendix records $57$ exceptional rows of the form
$n\; b_1:k_1\;\cdots\;b_t:k_t$.
The corresponding cubical sums are $c_j=k_j^3$
and the first boundary is $b_0=0$.
For every row, the boundaries increase,
$1\le L_j<R_j\le n+1$,
and the multiset equality in
Lemma~\ref{lem:balance} holds exactly.
These are $795$ simple integer-interval
identities. Exactly $10$ listed rows have
$n<173$ and exactly $47$ have $173\le n\le364$.
Together with the twenty-five algebraic-template lengths,
there are $22$ lower and $60$ upper seeds.

The following two lists give every omitted
base length without further matching data.
For any tabulated predecessor $m$ under
cube $c$, define $n=c-m-1$ and apply
Lemma~\ref{lem:suffix}. The lists are
ordered by cube, so every predecessor is
either a seed or was obtained in an earlier row.

\paragraph{First: 26 reflected lengths below 173.}
\begin{quote}
\footnotesize
\begin{verbatim}
cube     already constructed predecessors m
  8      0
 27      0,7
 64      0,7,19,26
125      0,7,19,26,37,44,56
216      44,56,63,66,68,80,82,85,87,
         100,103,105
\end{verbatim}
\end{quote}
Together with the $10$ exceptional lower seeds
and the twelve algebraic-template lower seeds,
these supply $48$ distinct lengths below 173.

\paragraph{Second: 49 reflected lengths in
$[173,255]\cup[339,364]$.}
\begin{quote}
\footnotesize
\begin{verbatim}
cube     previously constructed predecessors m
216      0,7,19,26,37
343      0,87,98,100,103,105,110,112,
         115,116,117,118,119,120,121,122,
         123,124,125,126,128,129,130,131,
         132,133,135,147,149,150,151,152,
         155,156,159
512      147,149,150,151,152,155,156,
         159,171
\end{verbatim}
\end{quote}
The $47$ upper exceptional rows, the thirteen
upper algebraic-template lengths and these $49$
reflections are exactly the $109$ integers
in $[173,255]\cup[339,364]$.
There is no circular dependency:
each predecessor $m$ is strictly smaller
than its constructed length $n$.

For example, $211=343-131-1$,
$340=512-171-1$, and $342=343-0-1$.
The first example replaces a previously
necessary-looking $53$-segment independent
row by one suffix reflection.

\begin{lemma}[Complete finite interval]\label{lem:base}
For every integer $173\le n\le364$,
a cube-sum permutation of $[1,n]$ exists.
\end{lemma}
\begin{proof}
The explicit arithmetic rows and the two
predecessor tables establish the claim for
$[173,255]\cup[339,364]$. For each of the
$83$ remaining lengths $256\le n\le338$,
take $m=511-n\in[173,255]$.
Then $m<n$ and $m+n+1=512=8^3$,
so Lemma~\ref{lem:suffix} applies.
\end{proof}

\subsection{Forced reflection peeling at the critical base}
\label{subsec:forced-reflection-peeling}

\begin{lemma}[Forced neighbor and reflected pair]
\label{lem:forced-neighbor}
Let $V$ be a finite set of positive integers, with
cube-sum adjacency $i\sim j$ whenever $i+j$ is a
cube. If $i\in V$ has exactly one neighbor $j\in V$,
every cube-sum permutation $p:V\to V$ satisfies
$p(i)=j$ and $p(j)=i$. Deleting $i,j$ (or only
$i$ when $i=j$) therefore preserves exactly the
number of cube-sum permutations of the induced
residual vertex set.
\end{lemma}
\begin{proof}
The $i$th row can use only image $j$.
Since the cube-sum graph is symmetric, column $i$
can only be supplied by row $j$. Both assignments
are therefore forced, and restriction or extension
gives mutually inverse maps between permutations.
\end{proof}

\begin{theorem}[Three forced bands]
\label{thm:forced-bands}
For every $173\le n\le181$, the following
three successive reflected bands are forced
in \emph{every} cube-sum permutation:
\[
\begin{array}{c|c|c}
\text{round}&\text{first endpoints }i& i+p(i)\\ \hline
1 &[125,342-n]&216\\
2 &[34,215-n]&64\\
3 &[95,276-n]&216
\end{array}
\]
Here each row is deleted together with its
reflected images before the next round.
Exactly $582-3n$ forced reflected pairs are
identified by these three bands.
\end{theorem}
\begin{proof}
Since $2n<512$, all relevant cube sums
are $8,27,64,125,216,343$. On the first
band, cubes up to $125$ give no positive
neighbor, and $343-i>n$; thus the unique
neighbor is $216-i$. Their image interval
is $[n-126,91]$.
On the second band, the possible $125-i$
is in this deleted image interval,
and $216-i>n$, leaving only $64-i$;
its images form $[n-151,30]$.
On the third band, the possible $125-i$
lies in the deleted second image interval,
and $343-i>n$, leaving $216-i$.
All partners are distinct and belong to
the still remaining vertices.
The band sizes are
$218-n$, $182-n$, and $182-n$, so their
sum is $582-3n$.
\end{proof}

\begin{theorem}[Exact counts at $173,174,175$]
\label{thm:forced-base}
For the cube-sum permutation counting function,
\[
\boxed{P(173)=1,\quad P(174)=1,\quad P(175)=3.}
\]
In particular the existence at $n=173$
is established without a determinant shape
or a precomputed endpoint-multiset identity.
\end{theorem}
\begin{proof}
At $n=173$, after the three forced bands
the additional batches of first endpoints
and their required cube sums are
\[
\begin{array}{c|l}
\text{round}& \text{first endpoints : cube}\\ \hline
4&\{4,5\}:8,\ [8,12]:27\\
5&\{45,46,106,107,108,122\}:216\\
6&\{31\}:64,\ \{172,173\}:216\\
7\!-\!15&(92,216),(1,8),(20,125),(111,125),
(13,125),\\
& (104,125),(6,8),(123,216),(32,64)
\end{array}
\]
Each first endpoint has exactly one neighbor
remaining when its round is reached;
this follows by direct subtraction from the
six cubes at most $343$. All $173$ vertices
are removed. Thus the resulting permutation
is unique, with fixed points $4,32,108$.

At $n=174$, the corresponding rounds are
\[
\begin{array}{c|l}
1\!-\!3&[125,168]:216,\ [34,41]:64,\ [95,102]:216\\
4&\{4\}:8,\ [8,11]:27\\
5&\{45,46,47,106,107,108\}:216\\
6&\{15\}:27,\ \{172,173,174\}:216\\
7\!-\!19&(113,216),(22,27),(3,125),(94,125),
(33,125),\\
&(124,125),(7,27),(105,216),(14,27),
(112,216),\\
&(21,27),(2,125),(93,125).
\end{array}
\]
Again every vertex is removed by forced choices,
so the permutation is unique, with fixed
points $4,108$.

At $n=175$ the first three bands are
$[125,167]$, $[34,40]$, $[95,101]$.
Use rounds $4$--$19$ of the preceding table,
but omit $\{4\}$ from round $4$ and
replace its interval $[8,11]$ by $[8,10]$.
The remaining ten vertices are precisely
\[
\{4,11,16,23,41,48,102,114,168,175\}.
\]
Their induced cube-sum graph is the loop at
$4$, the edge $(4,23)$, and the chordless
nine-cycle
\[
23-41-175-168-48-16-11-114-102-23.
\]
If $p(4)=4$, the odd cycle has exactly
two directed cycle covers, one for each
orientation. Otherwise $p(4)=23$
forces $p(23)=4$; the eight remaining
vertices form a path with a unique
cover by four reciprocal pairs. Thus
there are exactly $2+1=3$ permutations.
\end{proof}

\begin{corollary}[One fewer algebraic shape]
\label{cor:56-plus-forcing}
The proof of the finite interval $[173,364]$
may use the $56$ cubic-root/output-order
shapes in Appendix~\ref{app:reduced-shapes},
the six algebraic block families generating
$25$ further seed lengths, and the unique
forced matching of $173$ from
Theorem~\ref{thm:forced-base}.
Together these establish the same $82$
seed lengths as before. The excluded
$173$ determinant shape is redundant.
\end{corollary}

\subsection{The interval-overlap avalanche and its finite obstruction}
\label{subsec:interval-overlap-avalanche}

\begin{lemma}[Exact cubic reflection of a good interval]
\label{lem:reflection-overlap}
Suppose every integer $m\in[a,b]$, $0\le a\le b$,
has a legal cube-sum permutation. For a cube $c$,
every integer $n$ in the nonempty interval
\[
\bigl[
\max\{1,\lceil c/2\rceil,c-b-1\},\
c-a-1
\bigr]
\]
also has a legal cube-sum permutation.
\end{lemma}
\begin{proof}
Take $m=c-n-1$. By the displayed bounds,
$a\le m\le b$ and $0\le m<n$.
Keep the old permutation on $[1,m]$ and set
$p(i)=c-i$ on $[m+1,n]$. The suffix is mapped
bijectively to itself, with constant sum $c$.
\end{proof}

\begin{lemma}[Interval-overlap avalanche]
\label{lem:interval-avalanche}
Suppose every length in $[a,b]$ is legal,
where $1\le a\le b$. If a cube $c$ satisfies
\[
\boxed{a+b+2\le c\le2b+2},
\]
then every length in the larger interval
$[a,c-a-1]$ is legal.
\end{lemma}
\begin{proof}
By Lemma~\ref{lem:reflection-overlap},
the new interval of lengths includes $b+1$
because $a\le c-b-2\le b$ and
$c\le2(b+1)$. It continues through
$c-a-1\ge b+1$; it therefore meets the
already established interval immediately,
with no gap.
\end{proof}

\begin{theorem}[Infinite avalanche from the finite base]
\label{thm:interval-avalanche-tail}
If all lengths $173\le n\le364$ are legal,
then every length $n\ge173$ is legal,
without any additional finite certificate.
\end{theorem}
\begin{proof}
For $a=173$, $b=364$, the cube $9^3=729$
lies in the avalanche window $[539,730]$.
Consequently $[173,555]$ is legal.
Assume $[173,k^3-174]$ is legal for $k\ge9$.
The cube $(k+1)^3$ obeys
\[
k^3+1\le(k+1)^3\le2k^3-346,
\]
since the last inequality is equivalent to
\[
f(k)=k^3-3k^2-3k-347\ge0,
\]
with $f(9)=112$ and
$f(k+1)-f(k)=3k^2-3k-5>0$ for $k\ge9$.
Lemma~\ref{lem:interval-avalanche}
extends the good interval through
$(k+1)^3-174$. Induction on $k$ proves
the assertion.
\end{proof}

\paragraph{Why overlap does not by itself remove the finite base.}
At $n=65$, the domains $[1,63]$ of the
$64$-reflection and $[60,65]$ of the
$125$-reflection overlap and jointly cover
$[1,65]$. Yet
\[
S=\{3,40,59,64\},\qquad
N(S)=\{5,24,61\}
\]
under the available cube sums $8,27,64,125$,
so Hall's inequality fails.
Likewise no smaller legal length can produce
$n=173$ by a single suffix reflection:
the only possible predecessors are
$m=42$ and $m=169$, via cubes $216,343$.
For $m=42$, the six-row set
\[
[1,2]\cup[20,21]\cup[38,39]
\]
has just the four cube-sum neighbors
$[6,7]\cup[25,26]$.
For $m=169$, the fourteen-row set
\[
\begin{gathered}
[8,9]\cup[38,39]\cup[45,46]\cup[99,100]\\
{}\cup[129,130]\cup[136,137]\cup[160,161]
\end{gathered}
\]
has only the twelve neighbors
\[
\begin{gathered}
[18,19]\cup[25,26]\cup[55,56]\\
{}\cup[79,80]\cup[86,87]\cup[116,117].
\end{gathered}
\]
These are direct Hall obstructions.
There is also no cube in the interval
$[513,678]$, the necessary avalanche
window for extending $[173,338]$ through
$339$. Thus the avalanche lemma completely
removes the infinite tail's finite checks,
but \emph{does not derive the original
$173$--$364$ base interval without the
separately constructed finite witnesses}.
Such a claim would require an additional
local augmentation or gluing theorem.

\section{Existence for every length $n\ge173$}

\begin{theorem}
For every integer $n\ge173$, there is a
permutation $p$ of $[1,n]$ for which
every $i+p(i)$ is a perfect cube.
\end{theorem}
\begin{proof}
The finite interval $[173,364]$ is provided
by Lemma~\ref{lem:base}, from the
six explicit block families and $57$
algebraic cubic-order shapes.
The interval-overlap avalanche
Theorem~\ref{thm:interval-avalanche-tail}
then propagates it to every integer $n\ge173$.
No further finite construction is required.
\end{proof}

\paragraph{Exact logical status.}
The infinite propagation is purely symbolic,
including the full cube-gap inequality.
Each of the $57$ finite cubic-order shapes
is given by an integer Cramer-determinant
formula with a positive block-length vector,
not by an appeal to a matching solver.
The numerical root words, output orders,
and positivity checks for these $57$ shapes
are still indispensable \emph{finite case data}.
The exact audit program
\texttt{verify\_unimodular\_factory.py}
is a reproducibility aid, not a hidden
premise of the interval avalanche theorem.
The present argument does \emph{not}
establish a case-free existence proof
of the original base interval.

\clearpage
\appendix
\section{Redundant 57-row endpoint witnesses}
\label{app:table}
Every token $b:k$ means $p(i)=k^3-i$
up to endpoint $b$. For each row, verify
the same endpoint balance of
Lemma~\ref{lem:balance}.

\begin{multicols}{2}
\scriptsize\ttfamily
\setlength{\parindent}{0pt}
\setlength{\parskip}{2pt}
% 57 exceptional cube-sum interval identities, totaling 795 reflected blocks.
% The other 25 seeds follow from six algebraic permutation block-rearrangement families.
\noindent \textbf{82}: 2:4\allowbreak\ 5:2\allowbreak\ 21:3\allowbreak\ 42:4\allowbreak\ 61:5\allowbreak\ 63:4\allowbreak\ 82:5\par
\noindent \textbf{116}: 2:4\allowbreak\ 3:3\allowbreak\ 4:2\allowbreak\ 8:3\allowbreak\ 9:5\allowbreak\ 10:3\allowbreak\ 11:5\allowbreak\ 15:3\allowbreak\ 16:5\allowbreak\ 17:3\allowbreak\ 18:5\allowbreak\ 22:3\allowbreak\ 23:5\allowbreak\ 24:3\allowbreak\ 61:5\allowbreak\ 63:4\allowbreak\ 100:5\allowbreak\ 101:6\allowbreak\ 102:5\allowbreak\ 106:6\allowbreak\ 107:5\allowbreak\ 108:6\allowbreak\ 109:5\allowbreak\ 113:6\allowbreak\ 114:5\allowbreak\ 115:6\allowbreak\ 116:5\par
\noindent \textbf{118}: 1:3\allowbreak\ 2:2\allowbreak\ 3:3\allowbreak\ 4:2\allowbreak\ 5:3\allowbreak\ 6:2\allowbreak\ 7:5\allowbreak\ 8:3\allowbreak\ 9:5\allowbreak\ 10:3\allowbreak\ 11:5\allowbreak\ 12:3\allowbreak\ 14:5\allowbreak\ 15:3\allowbreak\ 16:5\allowbreak\ 17:3\allowbreak\ 18:5\allowbreak\ 19:3\allowbreak\ 21:5\allowbreak\ 22:3\allowbreak\ 23:5\allowbreak\ 24:3\allowbreak\ 25:5\allowbreak\ 26:3\allowbreak\ 98:5\allowbreak\ 99:6\allowbreak\ 100:5\allowbreak\ 101:6\allowbreak\ 102:5\allowbreak\ 103:6\allowbreak\ 105:5\allowbreak\ 106:6\allowbreak\ 107:5\allowbreak\ 108:6\allowbreak\ 109:5\allowbreak\ 110:6\allowbreak\ 112:5\allowbreak\ 113:6\allowbreak\ 114:5\allowbreak\ 115:6\allowbreak\ 116:5\allowbreak\ 117:6\allowbreak\ 118:5\par
\noindent \textbf{120}: 2:4\allowbreak\ 3:3\allowbreak\ 4:2\allowbreak\ 9:5\allowbreak\ 10:3\allowbreak\ 16:5\allowbreak\ 17:3\allowbreak\ 23:5\allowbreak\ 24:3\allowbreak\ 61:5\allowbreak\ 63:4\allowbreak\ 100:5\allowbreak\ 101:6\allowbreak\ 107:5\allowbreak\ 108:6\allowbreak\ 114:5\allowbreak\ 115:6\allowbreak\ 120:5\par
\noindent \textbf{121}: 2:4\allowbreak\ 3:3\allowbreak\ 9:5\allowbreak\ 10:3\allowbreak\ 16:5\allowbreak\ 17:3\allowbreak\ 23:5\allowbreak\ 24:3\allowbreak\ 61:5\allowbreak\ 63:4\allowbreak\ 100:5\allowbreak\ 101:6\allowbreak\ 107:5\allowbreak\ 108:6\allowbreak\ 114:5\allowbreak\ 115:6\allowbreak\ 121:5\par
\noindent \textbf{123}: 2:4\allowbreak\ 31:5\allowbreak\ 32:4\allowbreak\ 61:5\allowbreak\ 63:4\allowbreak\ 92:5\allowbreak\ 93:6\allowbreak\ 122:5\allowbreak\ 123:6\par
\noindent \textbf{125}: 3:5\allowbreak\ 4:2\allowbreak\ 29:5\allowbreak\ 30:4\allowbreak\ 33:5\allowbreak\ 34:4\allowbreak\ 90:5\allowbreak\ 91:6\allowbreak\ 94:5\allowbreak\ 95:6\allowbreak\ 120:5\allowbreak\ 121:6\allowbreak\ 124:5\allowbreak\ 125:6\par
\noindent \textbf{126}: 2:4\allowbreak\ 5:2\allowbreak\ 28:5\allowbreak\ 35:4\allowbreak\ 61:5\allowbreak\ 63:4\allowbreak\ 89:5\allowbreak\ 96:6\allowbreak\ 119:5\allowbreak\ 126:6\par
\noindent \textbf{127}: 1:5\allowbreak\ 4:2\allowbreak\ 5:3\allowbreak\ 6:2\allowbreak\ 11:5\allowbreak\ 12:3\allowbreak\ 18:5\allowbreak\ 19:3\allowbreak\ 25:5\allowbreak\ 26:3\allowbreak\ 27:5\allowbreak\ 30:4\allowbreak\ 31:5\allowbreak\ 36:4\allowbreak\ 88:5\allowbreak\ 91:6\allowbreak\ 92:5\allowbreak\ 97:6\allowbreak\ 102:5\allowbreak\ 103:6\allowbreak\ 109:5\allowbreak\ 110:6\allowbreak\ 116:5\allowbreak\ 117:6\allowbreak\ 118:5\allowbreak\ 121:6\allowbreak\ 122:5\allowbreak\ 127:6\par
\noindent \textbf{129}: 3:5\allowbreak\ 4:2\allowbreak\ 8:3\allowbreak\ 11:5\allowbreak\ 15:3\allowbreak\ 18:5\allowbreak\ 22:3\allowbreak\ 25:5\allowbreak\ 30:4\allowbreak\ 33:5\allowbreak\ 38:4\allowbreak\ 86:5\allowbreak\ 91:6\allowbreak\ 94:5\allowbreak\ 99:6\allowbreak\ 102:5\allowbreak\ 106:6\allowbreak\ 109:5\allowbreak\ 113:6\allowbreak\ 116:5\allowbreak\ 121:6\allowbreak\ 124:5\allowbreak\ 129:6\par
\noindent \textbf{173}: 7:2\allowbreak\ 12:3\allowbreak\ 14:5\allowbreak\ 19:3\allowbreak\ 21:5\allowbreak\ 42:4\allowbreak\ 103:6\allowbreak\ 105:5\allowbreak\ 110:6\allowbreak\ 112:5\allowbreak\ 173:6\par
\noindent \textbf{176}: 7:2\allowbreak\ 12:3\allowbreak\ 14:5\allowbreak\ 17:4\allowbreak\ 19:3\allowbreak\ 24:5\allowbreak\ 42:4\allowbreak\ 103:6\allowbreak\ 105:5\allowbreak\ 110:6\allowbreak\ 115:5\allowbreak\ 166:6\allowbreak\ 169:7\allowbreak\ 176:6\par
\noindent \textbf{177}: 3:5\allowbreak\ 4:2\allowbreak\ 8:3\allowbreak\ 11:5\allowbreak\ 15:3\allowbreak\ 18:5\allowbreak\ 22:3\allowbreak\ 25:5\allowbreak\ 30:4\allowbreak\ 33:5\allowbreak\ 38:4\allowbreak\ 91:6\allowbreak\ 94:5\allowbreak\ 99:6\allowbreak\ 102:5\allowbreak\ 106:6\allowbreak\ 109:5\allowbreak\ 113:6\allowbreak\ 116:5\allowbreak\ 121:6\allowbreak\ 124:5\allowbreak\ 177:6\par
\noindent \textbf{179}: 3:5\allowbreak\ 11:3\allowbreak\ 15:4\allowbreak\ 23:3\allowbreak\ 30:4\allowbreak\ 33:5\allowbreak\ 40:4\allowbreak\ 48:6\allowbreak\ 52:4\allowbreak\ 91:6\allowbreak\ 94:5\allowbreak\ 121:6\allowbreak\ 124:5\allowbreak\ 163:6\allowbreak\ 167:7\allowbreak\ 175:6\allowbreak\ 179:7\par
\noindent \textbf{180}: 5:3\allowbreak\ 7:5\allowbreak\ 12:3\allowbreak\ 14:5\allowbreak\ 19:3\allowbreak\ 21:5\allowbreak\ 26:3\allowbreak\ 28:5\allowbreak\ 35:4\allowbreak\ 96:6\allowbreak\ 98:5\allowbreak\ 103:6\allowbreak\ 105:5\allowbreak\ 110:6\allowbreak\ 112:5\allowbreak\ 117:6\allowbreak\ 119:5\allowbreak\ 180:6\par
\noindent \textbf{181}: 3:5\allowbreak\ 4:2\allowbreak\ 29:5\allowbreak\ 30:4\allowbreak\ 33:5\allowbreak\ 34:4\allowbreak\ 91:6\allowbreak\ 94:5\allowbreak\ 95:6\allowbreak\ 120:5\allowbreak\ 121:6\allowbreak\ 124:5\allowbreak\ 181:6\par
\noindent \textbf{184}: 7:2\allowbreak\ 31:4\allowbreak\ 32:6\allowbreak\ 56:4\allowbreak\ 159:6\allowbreak\ 183:7\allowbreak\ 184:6\par
\noindent \textbf{185}: 21:5\allowbreak\ 26:3\allowbreak\ 30:4\allowbreak\ 51:6\allowbreak\ 58:4\allowbreak\ 91:6\allowbreak\ 112:5\allowbreak\ 178:6\allowbreak\ 185:7\par
\noindent \textbf{188}: 2:5\allowbreak\ 5:4\allowbreak\ 7:2\allowbreak\ 30:4\allowbreak\ 32:6\allowbreak\ 61:4\allowbreak\ 91:6\allowbreak\ 93:5\allowbreak\ 154:6\allowbreak\ 157:7\allowbreak\ 159:6\allowbreak\ 188:7\par
\noindent \textbf{194}: 7:2\allowbreak\ 12:3\allowbreak\ 14:5\allowbreak\ 19:3\allowbreak\ 21:5\allowbreak\ 103:6\allowbreak\ 105:5\allowbreak\ 110:6\allowbreak\ 112:5\allowbreak\ 194:6\par
\noindent \textbf{197}: 8:5\allowbreak\ 18:3\allowbreak\ 55:6\allowbreak\ 63:4\allowbreak\ 91:6\allowbreak\ 99:5\allowbreak\ 182:6\allowbreak\ 190:7\allowbreak\ 197:6\par
\noindent \textbf{198}: 2:4\allowbreak\ 5:2\allowbreak\ 28:5\allowbreak\ 35:4\allowbreak\ 53:6\allowbreak\ 61:5\allowbreak\ 63:4\allowbreak\ 71:5\allowbreak\ 96:6\allowbreak\ 119:5\allowbreak\ 144:6\allowbreak\ 162:7\allowbreak\ 180:6\allowbreak\ 198:7\par
\noindent \textbf{199}: 2:4\allowbreak\ 5:2\allowbreak\ 16:4\allowbreak\ 22:6\allowbreak\ 41:4\allowbreak\ 47:6\allowbreak\ 58:4\allowbreak\ 61:5\allowbreak\ 63:4\allowbreak\ 66:5\allowbreak\ 149:6\allowbreak\ 168:7\allowbreak\ 174:6\allowbreak\ 193:7\allowbreak\ 199:6\par
\noindent \textbf{200}: 26:3\allowbreak\ 37:4\allowbreak\ 51:6\allowbreak\ 73:5\allowbreak\ 142:6\allowbreak\ 164:7\allowbreak\ 178:6\allowbreak\ 200:7\par
\noindent \textbf{201}: 1:3\allowbreak\ 12:4\allowbreak\ 25:3\allowbreak\ 37:4\allowbreak\ 50:6\allowbreak\ 62:5\allowbreak\ 63:4\allowbreak\ 141:6\allowbreak\ 164:7\allowbreak\ 177:6\allowbreak\ 190:7\allowbreak\ 201:6\par
\noindent \textbf{202}: 13:4\allowbreak\ 25:6\allowbreak\ 38:4\allowbreak\ 50:6\allowbreak\ 63:4\allowbreak\ 152:6\allowbreak\ 165:7\allowbreak\ 177:6\allowbreak\ 190:7\allowbreak\ 202:6\par
\noindent \textbf{203}: 12:4\allowbreak\ 25:6\allowbreak\ 26:3\allowbreak\ 37:4\allowbreak\ 51:6\allowbreak\ 62:4\allowbreak\ 152:6\allowbreak\ 164:7\allowbreak\ 178:6\allowbreak\ 189:7\allowbreak\ 203:6\par
\noindent \textbf{204}: 11:4\allowbreak\ 15:6\allowbreak\ 26:3\allowbreak\ 36:4\allowbreak\ 51:6\allowbreak\ 61:5\allowbreak\ 142:6\allowbreak\ 163:7\allowbreak\ 178:6\allowbreak\ 188:7\allowbreak\ 204:6\par
\noindent \textbf{205}: 2:4\allowbreak\ 5:2\allowbreak\ 21:3\allowbreak\ 42:4\allowbreak\ 46:6\allowbreak\ 61:5\allowbreak\ 63:4\allowbreak\ 78:5\allowbreak\ 137:6\allowbreak\ 169:7\allowbreak\ 173:6\allowbreak\ 205:7\par
\noindent \textbf{206}: 7:2\allowbreak\ 18:3\allowbreak\ 19:4\allowbreak\ 20:6\allowbreak\ 44:4\allowbreak\ 45:6\allowbreak\ 55:5\allowbreak\ 56:4\allowbreak\ 79:5\allowbreak\ 136:6\allowbreak\ 146:7\allowbreak\ 147:6\allowbreak\ 171:7\allowbreak\ 172:6\allowbreak\ 206:7\par
\noindent \textbf{233}: 7:2\allowbreak\ 9:6\allowbreak\ 18:4\allowbreak\ 43:5\allowbreak\ 45:6\allowbreak\ 70:5\allowbreak\ 81:6\allowbreak\ 106:5\allowbreak\ 109:6\allowbreak\ 134:7\allowbreak\ 136:6\allowbreak\ 170:7\allowbreak\ 172:6\allowbreak\ 197:7\allowbreak\ 208:6\allowbreak\ 233:7\par
\noindent \textbf{234}: 16:6\allowbreak\ 36:5\allowbreak\ 52:6\allowbreak\ 72:5\allowbreak\ 88:6\allowbreak\ 108:5\allowbreak\ 127:7\allowbreak\ 143:6\allowbreak\ 163:7\allowbreak\ 179:6\allowbreak\ 199:7\allowbreak\ 215:6\allowbreak\ 234:7\par
\noindent \textbf{235}: 7:2\allowbreak\ 9:6\allowbreak\ 17:3\allowbreak\ 43:5\allowbreak\ 45:6\allowbreak\ 79:5\allowbreak\ 81:6\allowbreak\ 107:5\allowbreak\ 134:7\allowbreak\ 136:6\allowbreak\ 170:7\allowbreak\ 172:6\allowbreak\ 206:7\allowbreak\ 208:6\allowbreak\ 235:7\par
\noindent \textbf{236}: 7:2\allowbreak\ 18:6\allowbreak\ 43:5\allowbreak\ 45:6\allowbreak\ 54:4\allowbreak\ 79:5\allowbreak\ 81:6\allowbreak\ 106:5\allowbreak\ 134:7\allowbreak\ 145:6\allowbreak\ 170:7\allowbreak\ 172:6\allowbreak\ 206:7\allowbreak\ 208:6\allowbreak\ 236:7\par
\noindent \textbf{243}: 25:4\allowbreak\ 61:5\allowbreak\ 63:4\allowbreak\ 86:6\allowbreak\ 99:5\allowbreak\ 190:7\allowbreak\ 213:6\allowbreak\ 243:7\par
\noindent \textbf{245}: 27:4\allowbreak\ 61:5\allowbreak\ 88:6\allowbreak\ 97:5\allowbreak\ 188:7\allowbreak\ 215:6\allowbreak\ 245:7\par
\noindent \textbf{246}: 2:6\allowbreak\ 24:3\allowbreak\ 28:6\allowbreak\ 36:5\allowbreak\ 38:4\allowbreak\ 60:5\allowbreak\ 61:6\allowbreak\ 63:4\allowbreak\ 64:6\allowbreak\ 96:5\allowbreak\ 127:7\allowbreak\ 129:6\allowbreak\ 151:7\allowbreak\ 155:6\allowbreak\ 187:7\allowbreak\ 188:6\allowbreak\ 190:7\allowbreak\ 191:6\allowbreak\ 246:7\par
\noindent \textbf{247}: 2:6\allowbreak\ 3:2\allowbreak\ 4:3\allowbreak\ 5:2\allowbreak\ 22:6\allowbreak\ 23:3\allowbreak\ 29:6\allowbreak\ 36:5\allowbreak\ 38:4\allowbreak\ 41:5\allowbreak\ 58:4\allowbreak\ 59:5\allowbreak\ 61:6\allowbreak\ 63:4\allowbreak\ 65:6\allowbreak\ 95:5\allowbreak\ 127:7\allowbreak\ 129:6\allowbreak\ 132:7\allowbreak\ 149:6\allowbreak\ 150:7\allowbreak\ 156:6\allowbreak\ 186:7\allowbreak\ 188:6\allowbreak\ 190:7\allowbreak\ 192:6\allowbreak\ 247:7\par
\noindent \textbf{248}: 2:6\allowbreak\ 5:2\allowbreak\ 30:6\allowbreak\ 36:5\allowbreak\ 38:4\allowbreak\ 41:5\allowbreak\ 58:4\allowbreak\ 61:6\allowbreak\ 63:4\allowbreak\ 66:6\allowbreak\ 94:5\allowbreak\ 127:7\allowbreak\ 129:6\allowbreak\ 132:7\allowbreak\ 157:6\allowbreak\ 185:7\allowbreak\ 188:6\allowbreak\ 190:7\allowbreak\ 193:6\allowbreak\ 248:7\par
\noindent \textbf{249}: 1:3\allowbreak\ 6:2\allowbreak\ 31:6\allowbreak\ 42:5\allowbreak\ 57:4\allowbreak\ 62:6\allowbreak\ 63:4\allowbreak\ 67:6\allowbreak\ 93:5\allowbreak\ 133:7\allowbreak\ 158:6\allowbreak\ 184:7\allowbreak\ 189:6\allowbreak\ 190:7\allowbreak\ 194:6\allowbreak\ 249:7\par
\noindent \textbf{250}: 7:2\allowbreak\ 32:6\allowbreak\ 43:5\allowbreak\ 56:4\allowbreak\ 68:6\allowbreak\ 92:5\allowbreak\ 134:7\allowbreak\ 159:6\allowbreak\ 183:7\allowbreak\ 195:6\allowbreak\ 250:7\par
\noindent \textbf{251}: 7:2\allowbreak\ 19:3\allowbreak\ 33:6\allowbreak\ 55:5\allowbreak\ 69:6\allowbreak\ 91:5\allowbreak\ 146:7\allowbreak\ 160:6\allowbreak\ 182:7\allowbreak\ 196:6\allowbreak\ 251:7\par
\noindent \textbf{252}: 7:2\allowbreak\ 9:6\allowbreak\ 18:4\allowbreak\ 34:6\allowbreak\ 43:5\allowbreak\ 45:6\allowbreak\ 54:5\allowbreak\ 81:6\allowbreak\ 90:5\allowbreak\ 134:7\allowbreak\ 136:6\allowbreak\ 145:7\allowbreak\ 161:6\allowbreak\ 170:7\allowbreak\ 172:6\allowbreak\ 181:7\allowbreak\ 208:6\allowbreak\ 252:7\par
\noindent \textbf{253}: 2:4\allowbreak\ 24:3\allowbreak\ 25:6\allowbreak\ 27:4\allowbreak\ 28:6\allowbreak\ 35:4\allowbreak\ 60:5\allowbreak\ 64:6\allowbreak\ 86:5\allowbreak\ 88:6\allowbreak\ 89:5\allowbreak\ 151:7\allowbreak\ 152:6\allowbreak\ 154:7\allowbreak\ 155:6\allowbreak\ 187:7\allowbreak\ 191:6\allowbreak\ 213:7\allowbreak\ 215:6\allowbreak\ 253:7\par
\noindent \textbf{339}: 6:3\allowbreak\ 20:6\allowbreak\ 23:4\allowbreak\ 26:3\allowbreak\ 172:6\allowbreak\ 302:8\allowbreak\ 316:7\allowbreak\ 322:8\allowbreak\ 339:7\par
\noindent \textbf{341}: 1:3\allowbreak\ 21:7\allowbreak\ 62:6\allowbreak\ 63:4\allowbreak\ 189:6\allowbreak\ 190:7\allowbreak\ 317:8\allowbreak\ 341:7\par
\noindent \textbf{343}: 16:5\allowbreak\ 30:4\allowbreak\ 46:6\allowbreak\ 47:4\allowbreak\ 91:6\allowbreak\ 107:5\allowbreak\ 168:6\allowbreak\ 326:8\allowbreak\ 342:7\allowbreak\ 343:8\par
\noindent \textbf{344}: 40:7\allowbreak\ 42:6\allowbreak\ 82:5\allowbreak\ 133:6\allowbreak\ 167:7\allowbreak\ 169:8\allowbreak\ 173:7\allowbreak\ 175:6\allowbreak\ 209:7\allowbreak\ 302:8\allowbreak\ 342:7\allowbreak\ 344:8\par
\noindent \textbf{346}: 8:5\allowbreak\ 49:7\allowbreak\ 63:4\allowbreak\ 74:6\allowbreak\ 99:5\allowbreak\ 165:6\allowbreak\ 177:8\allowbreak\ 190:7\allowbreak\ 201:6\allowbreak\ 218:7\allowbreak\ 346:8\par
\noindent \textbf{347}: 9:3\allowbreak\ 37:4\allowbreak\ 46:6\allowbreak\ 51:4\allowbreak\ 161:6\allowbreak\ 164:7\allowbreak\ 330:8\allowbreak\ 342:7\allowbreak\ 347:8\par
\noindent \textbf{348}: 10:6\allowbreak\ 36:5\allowbreak\ 46:6\allowbreak\ 52:5\allowbreak\ 101:6\allowbreak\ 127:7\allowbreak\ 137:6\allowbreak\ 163:7\allowbreak\ 270:8\allowbreak\ 342:7\allowbreak\ 348:8\par
\noindent \textbf{349}: 46:7\allowbreak\ 53:4\allowbreak\ 71:6\allowbreak\ 78:5\allowbreak\ 162:6\allowbreak\ 169:8\allowbreak\ 198:6\allowbreak\ 205:7\allowbreak\ 215:6\allowbreak\ 349:8\par
\noindent \textbf{350}: 9:7\allowbreak\ 17:4\allowbreak\ 26:3\allowbreak\ 54:4\allowbreak\ 161:6\allowbreak\ 169:8\allowbreak\ 178:6\allowbreak\ 350:8\par
\noindent \textbf{353}: 6:3\allowbreak\ 7:2\allowbreak\ 20:6\allowbreak\ 25:3\allowbreak\ 43:5\allowbreak\ 57:4\allowbreak\ 116:6\allowbreak\ 134:7\allowbreak\ 158:6\allowbreak\ 285:8\allowbreak\ 316:7\allowbreak\ 353:8\par
\noindent \textbf{354}: 5:7\allowbreak\ 21:4\allowbreak\ 22:3\allowbreak\ 58:4\allowbreak\ 157:6\allowbreak\ 169:8\allowbreak\ 173:7\allowbreak\ 174:6\allowbreak\ 338:8\allowbreak\ 342:7\allowbreak\ 354:8\par
\noindent \textbf{358}: 6:3\allowbreak\ 7:2\allowbreak\ 43:7\allowbreak\ 62:4\allowbreak\ 85:5\allowbreak\ 153:6\allowbreak\ 176:8\allowbreak\ 189:6\allowbreak\ 212:7\allowbreak\ 358:8\par
\noindent \textbf{363}: 4:7\allowbreak\ 6:6\allowbreak\ 21:4\allowbreak\ 42:6\allowbreak\ 67:5\allowbreak\ 133:6\allowbreak\ 148:7\allowbreak\ 169:8\allowbreak\ 173:7\allowbreak\ 300:8\allowbreak\ 342:7\allowbreak\ 363:8\par

\end{multicols}


\clearpage
\section{The 57 algebraic cubic-order shapes}
\label{app:algebraic-shapes}
For each target $n$, the line gives a cube-root
word $k_1,\ldots,k_t$ and an output-block order
$\pi$. Its block lengths are not assumed or
listed: they are the integer determinants
in Theorem~\ref{thm:unimodular-block-factory}.
The target length is a checksum for their sum.
All listed lengths are strictly positive.
\begin{multicols}{2}
\scriptsize
\setlength{\parindent}{0pt}
\setlength{\parskip}{2pt}
% 57 finite algebraic shapes: each line specifies cubic roots and output order, no endpoint table.
% n is only a checksum: the determinant length formulas derive n as their sum.
% Read alongside UNIMODULAR_BLOCK_FACTORY.md and verify_unimodular_factory.py.
\noindent \textbf{82}: roots \texttt{4,\allowbreak 2,\allowbreak 3,\allowbreak 4,\allowbreak 5,\allowbreak 4,\allowbreak 5}; order \texttt{6,\allowbreak 2,\allowbreak 3,\allowbreak 4,\allowbreak 7,\allowbreak 1,\allowbreak 5}\par
\noindent \textbf{116}: roots \texttt{4,\allowbreak 3,\allowbreak 2,\allowbreak 3,\allowbreak 5,\allowbreak 3,\allowbreak 5,\allowbreak 3,\allowbreak 5,\allowbreak 3,\allowbreak 5,\allowbreak 3,\allowbreak 5,\allowbreak 3,\allowbreak 5,\allowbreak 4,\allowbreak 5,\allowbreak 6,\allowbreak 5,\allowbreak 6,\allowbreak 5,\allowbreak 6,\allowbreak 5,\allowbreak 6,\allowbreak 5,\allowbreak 6,\allowbreak 5}; order \texttt{16,\allowbreak 14,\allowbreak 3,\allowbreak 12,\allowbreak 27,\allowbreak 10,\allowbreak 25,\allowbreak 8,\allowbreak 23,\allowbreak 6,\allowbreak 21,\allowbreak 4,\allowbreak 19,\allowbreak 2,\allowbreak 17,\allowbreak 1,\allowbreak 15,\allowbreak 26,\allowbreak 13,\allowbreak 24,\allowbreak 11,\allowbreak 22,\allowbreak 9,\allowbreak 20,\allowbreak 7,\allowbreak 18,\allowbreak 5}\par
\noindent \textbf{118}: roots \texttt{3,\allowbreak 2,\allowbreak 3,\allowbreak 2,\allowbreak 3,\allowbreak 2,\allowbreak 5,\allowbreak 3,\allowbreak 5,\allowbreak 3,\allowbreak 5,\allowbreak 3,\allowbreak 5,\allowbreak 3,\allowbreak 5,\allowbreak 3,\allowbreak 5,\allowbreak 3,\allowbreak 5,\allowbreak 3,\allowbreak 5,\allowbreak 3,\allowbreak 5,\allowbreak 3,\allowbreak 5,\allowbreak 6,\allowbreak 5,\allowbreak 6,\allowbreak 5,\allowbreak 6,\allowbreak 5,\allowbreak 6,\allowbreak 5,\allowbreak 6,\allowbreak 5,\allowbreak 6,\allowbreak 5,\allowbreak 6,\allowbreak 5,\allowbreak 6,\allowbreak 5,\allowbreak 6,\allowbreak 5}; order \texttt{24,\allowbreak 6,\allowbreak 22,\allowbreak 4,\allowbreak 20,\allowbreak 2,\allowbreak 43,\allowbreak 18,\allowbreak 41,\allowbreak 16,\allowbreak 39,\allowbreak 14,\allowbreak 37,\allowbreak 12,\allowbreak 35,\allowbreak 10,\allowbreak 33,\allowbreak 8,\allowbreak 31,\allowbreak 5,\allowbreak 29,\allowbreak 3,\allowbreak 27,\allowbreak 1,\allowbreak 25,\allowbreak 42,\allowbreak 23,\allowbreak 40,\allowbreak 21,\allowbreak 38,\allowbreak 19,\allowbreak 36,\allowbreak 17,\allowbreak 34,\allowbreak 15,\allowbreak 32,\allowbreak 13,\allowbreak 30,\allowbreak 11,\allowbreak 28,\allowbreak 9,\allowbreak 26,\allowbreak 7}\par
\noindent \textbf{120}: roots \texttt{4,\allowbreak 3,\allowbreak 2,\allowbreak 5,\allowbreak 3,\allowbreak 5,\allowbreak 3,\allowbreak 5,\allowbreak 3,\allowbreak 5,\allowbreak 4,\allowbreak 5,\allowbreak 6,\allowbreak 5,\allowbreak 6,\allowbreak 5,\allowbreak 6,\allowbreak 5}; order \texttt{11,\allowbreak 9,\allowbreak 3,\allowbreak 18,\allowbreak 7,\allowbreak 16,\allowbreak 5,\allowbreak 14,\allowbreak 2,\allowbreak 12,\allowbreak 1,\allowbreak 10,\allowbreak 17,\allowbreak 8,\allowbreak 15,\allowbreak 6,\allowbreak 13,\allowbreak 4}\par
\noindent \textbf{121}: roots \texttt{4,\allowbreak 3,\allowbreak 5,\allowbreak 3,\allowbreak 5,\allowbreak 3,\allowbreak 5,\allowbreak 3,\allowbreak 5,\allowbreak 4,\allowbreak 5,\allowbreak 6,\allowbreak 5,\allowbreak 6,\allowbreak 5,\allowbreak 6,\allowbreak 5}; order \texttt{10,\allowbreak 8,\allowbreak 17,\allowbreak 6,\allowbreak 15,\allowbreak 4,\allowbreak 13,\allowbreak 2,\allowbreak 11,\allowbreak 1,\allowbreak 9,\allowbreak 16,\allowbreak 7,\allowbreak 14,\allowbreak 5,\allowbreak 12,\allowbreak 3}\par
\noindent \textbf{123}: roots \texttt{4,\allowbreak 5,\allowbreak 4,\allowbreak 5,\allowbreak 4,\allowbreak 5,\allowbreak 6,\allowbreak 5,\allowbreak 6}; order \texttt{5,\allowbreak 8,\allowbreak 3,\allowbreak 6,\allowbreak 1,\allowbreak 4,\allowbreak 9,\allowbreak 2,\allowbreak 7}\par
\noindent \textbf{125}: roots \texttt{5,\allowbreak 2,\allowbreak 5,\allowbreak 4,\allowbreak 5,\allowbreak 4,\allowbreak 5,\allowbreak 6,\allowbreak 5,\allowbreak 6,\allowbreak 5,\allowbreak 6,\allowbreak 5,\allowbreak 6}; order \texttt{13,\allowbreak 2,\allowbreak 11,\allowbreak 6,\allowbreak 9,\allowbreak 4,\allowbreak 7,\allowbreak 14,\allowbreak 5,\allowbreak 12,\allowbreak 3,\allowbreak 10,\allowbreak 1,\allowbreak 8}\par
\noindent \textbf{126}: roots \texttt{4,\allowbreak 2,\allowbreak 5,\allowbreak 4,\allowbreak 5,\allowbreak 4,\allowbreak 5,\allowbreak 6,\allowbreak 5,\allowbreak 6}; order \texttt{6,\allowbreak 2,\allowbreak 9,\allowbreak 4,\allowbreak 7,\allowbreak 1,\allowbreak 5,\allowbreak 10,\allowbreak 3,\allowbreak 8}\par
\noindent \textbf{127}: roots \texttt{5,\allowbreak 2,\allowbreak 3,\allowbreak 2,\allowbreak 5,\allowbreak 3,\allowbreak 5,\allowbreak 3,\allowbreak 5,\allowbreak 3,\allowbreak 5,\allowbreak 4,\allowbreak 5,\allowbreak 4,\allowbreak 5,\allowbreak 6,\allowbreak 5,\allowbreak 6,\allowbreak 5,\allowbreak 6,\allowbreak 5,\allowbreak 6,\allowbreak 5,\allowbreak 6,\allowbreak 5,\allowbreak 6,\allowbreak 5,\allowbreak 6}; order \texttt{10,\allowbreak 4,\allowbreak 27,\allowbreak 2,\allowbreak 25,\allowbreak 8,\allowbreak 23,\allowbreak 6,\allowbreak 21,\allowbreak 3,\allowbreak 19,\allowbreak 14,\allowbreak 17,\allowbreak 12,\allowbreak 15,\allowbreak 28,\allowbreak 13,\allowbreak 26,\allowbreak 11,\allowbreak 24,\allowbreak 9,\allowbreak 22,\allowbreak 7,\allowbreak 20,\allowbreak 5,\allowbreak 18,\allowbreak 1,\allowbreak 16}\par
\noindent \textbf{129}: roots \texttt{5,\allowbreak 2,\allowbreak 3,\allowbreak 5,\allowbreak 3,\allowbreak 5,\allowbreak 3,\allowbreak 5,\allowbreak 4,\allowbreak 5,\allowbreak 4,\allowbreak 5,\allowbreak 6,\allowbreak 5,\allowbreak 6,\allowbreak 5,\allowbreak 6,\allowbreak 5,\allowbreak 6,\allowbreak 5,\allowbreak 6,\allowbreak 5,\allowbreak 6}; order \texttt{22,\allowbreak 2,\allowbreak 7,\allowbreak 20,\allowbreak 5,\allowbreak 18,\allowbreak 3,\allowbreak 16,\allowbreak 11,\allowbreak 14,\allowbreak 9,\allowbreak 12,\allowbreak 23,\allowbreak 10,\allowbreak 21,\allowbreak 8,\allowbreak 19,\allowbreak 6,\allowbreak 17,\allowbreak 4,\allowbreak 15,\allowbreak 1,\allowbreak 13}\par
\noindent \textbf{173}: roots \texttt{2,\allowbreak 3,\allowbreak 5,\allowbreak 3,\allowbreak 5,\allowbreak 4,\allowbreak 6,\allowbreak 5,\allowbreak 6,\allowbreak 5,\allowbreak 6}; order \texttt{1,\allowbreak 4,\allowbreak 10,\allowbreak 2,\allowbreak 8,\allowbreak 6,\allowbreak 11,\allowbreak 5,\allowbreak 9,\allowbreak 3,\allowbreak 7}\par
\noindent \textbf{176}: roots \texttt{2,\allowbreak 3,\allowbreak 5,\allowbreak 4,\allowbreak 3,\allowbreak 5,\allowbreak 4,\allowbreak 6,\allowbreak 5,\allowbreak 6,\allowbreak 5,\allowbreak 6,\allowbreak 7,\allowbreak 6}; order \texttt{1,\allowbreak 5,\allowbreak 11,\allowbreak 2,\allowbreak 9,\allowbreak 7,\allowbreak 14,\allowbreak 4,\allowbreak 12,\allowbreak 6,\allowbreak 10,\allowbreak 3,\allowbreak 8,\allowbreak 13}\par
\noindent \textbf{177}: roots \texttt{5,\allowbreak 2,\allowbreak 3,\allowbreak 5,\allowbreak 3,\allowbreak 5,\allowbreak 3,\allowbreak 5,\allowbreak 4,\allowbreak 5,\allowbreak 4,\allowbreak 6,\allowbreak 5,\allowbreak 6,\allowbreak 5,\allowbreak 6,\allowbreak 5,\allowbreak 6,\allowbreak 5,\allowbreak 6,\allowbreak 5,\allowbreak 6}; order \texttt{21,\allowbreak 2,\allowbreak 7,\allowbreak 19,\allowbreak 5,\allowbreak 17,\allowbreak 3,\allowbreak 15,\allowbreak 11,\allowbreak 13,\allowbreak 9,\allowbreak 22,\allowbreak 10,\allowbreak 20,\allowbreak 8,\allowbreak 18,\allowbreak 6,\allowbreak 16,\allowbreak 4,\allowbreak 14,\allowbreak 1,\allowbreak 12}\par
\noindent \textbf{179}: roots \texttt{5,\allowbreak 3,\allowbreak 4,\allowbreak 3,\allowbreak 4,\allowbreak 5,\allowbreak 4,\allowbreak 6,\allowbreak 4,\allowbreak 6,\allowbreak 5,\allowbreak 6,\allowbreak 5,\allowbreak 6,\allowbreak 7,\allowbreak 6,\allowbreak 7}; order \texttt{13,\allowbreak 4,\allowbreak 9,\allowbreak 2,\allowbreak 7,\allowbreak 11,\allowbreak 5,\allowbreak 16,\allowbreak 3,\allowbreak 14,\allowbreak 6,\allowbreak 12,\allowbreak 1,\allowbreak 10,\allowbreak 17,\allowbreak 8,\allowbreak 15}\par
\noindent \textbf{180}: roots \texttt{3,\allowbreak 5,\allowbreak 3,\allowbreak 5,\allowbreak 3,\allowbreak 5,\allowbreak 3,\allowbreak 5,\allowbreak 4,\allowbreak 6,\allowbreak 5,\allowbreak 6,\allowbreak 5,\allowbreak 6,\allowbreak 5,\allowbreak 6,\allowbreak 5,\allowbreak 6}; order \texttt{7,\allowbreak 17,\allowbreak 5,\allowbreak 15,\allowbreak 3,\allowbreak 13,\allowbreak 1,\allowbreak 11,\allowbreak 9,\allowbreak 18,\allowbreak 8,\allowbreak 16,\allowbreak 6,\allowbreak 14,\allowbreak 4,\allowbreak 12,\allowbreak 2,\allowbreak 10}\par
\noindent \textbf{181}: roots \texttt{5,\allowbreak 2,\allowbreak 5,\allowbreak 4,\allowbreak 5,\allowbreak 4,\allowbreak 6,\allowbreak 5,\allowbreak 6,\allowbreak 5,\allowbreak 6,\allowbreak 5,\allowbreak 6}; order \texttt{12,\allowbreak 2,\allowbreak 10,\allowbreak 6,\allowbreak 8,\allowbreak 4,\allowbreak 13,\allowbreak 5,\allowbreak 11,\allowbreak 3,\allowbreak 9,\allowbreak 1,\allowbreak 7}\par
\noindent \textbf{184}: roots \texttt{2,\allowbreak 4,\allowbreak 6,\allowbreak 4,\allowbreak 6,\allowbreak 7,\allowbreak 6}; order \texttt{1,\allowbreak 4,\allowbreak 7,\allowbreak 2,\allowbreak 5,\allowbreak 6,\allowbreak 3}\par
\noindent \textbf{185}: roots \texttt{5,\allowbreak 3,\allowbreak 4,\allowbreak 6,\allowbreak 4,\allowbreak 6,\allowbreak 5,\allowbreak 6,\allowbreak 7}; order \texttt{2,\allowbreak 5,\allowbreak 7,\allowbreak 3,\allowbreak 8,\allowbreak 1,\allowbreak 6,\allowbreak 9,\allowbreak 4}\par
\noindent \textbf{188}: roots \texttt{5,\allowbreak 4,\allowbreak 2,\allowbreak 4,\allowbreak 6,\allowbreak 4,\allowbreak 6,\allowbreak 5,\allowbreak 6,\allowbreak 7,\allowbreak 6,\allowbreak 7}; order \texttt{3,\allowbreak 6,\allowbreak 8,\allowbreak 4,\allowbreak 11,\allowbreak 2,\allowbreak 9,\allowbreak 1,\allowbreak 7,\allowbreak 12,\allowbreak 5,\allowbreak 10}\par
\noindent \textbf{194}: roots \texttt{2,\allowbreak 3,\allowbreak 5,\allowbreak 3,\allowbreak 5,\allowbreak 6,\allowbreak 5,\allowbreak 6,\allowbreak 5,\allowbreak 6}; order \texttt{1,\allowbreak 4,\allowbreak 9,\allowbreak 2,\allowbreak 7,\allowbreak 10,\allowbreak 5,\allowbreak 8,\allowbreak 3,\allowbreak 6}\par
\noindent \textbf{197}: roots \texttt{5,\allowbreak 3,\allowbreak 6,\allowbreak 4,\allowbreak 6,\allowbreak 5,\allowbreak 6,\allowbreak 7,\allowbreak 6}; order \texttt{4,\allowbreak 2,\allowbreak 9,\allowbreak 6,\allowbreak 7,\allowbreak 1,\allowbreak 5,\allowbreak 8,\allowbreak 3}\par
\noindent \textbf{198}: roots \texttt{4,\allowbreak 2,\allowbreak 5,\allowbreak 4,\allowbreak 6,\allowbreak 5,\allowbreak 4,\allowbreak 5,\allowbreak 6,\allowbreak 5,\allowbreak 6,\allowbreak 7,\allowbreak 6,\allowbreak 7}; order \texttt{7,\allowbreak 2,\allowbreak 10,\allowbreak 4,\allowbreak 13,\allowbreak 8,\allowbreak 1,\allowbreak 6,\allowbreak 11,\allowbreak 3,\allowbreak 9,\allowbreak 14,\allowbreak 5,\allowbreak 12}\par
\noindent \textbf{199}: roots \texttt{4,\allowbreak 2,\allowbreak 4,\allowbreak 6,\allowbreak 4,\allowbreak 6,\allowbreak 4,\allowbreak 5,\allowbreak 4,\allowbreak 5,\allowbreak 6,\allowbreak 7,\allowbreak 6,\allowbreak 7,\allowbreak 6}; order \texttt{9,\allowbreak 2,\allowbreak 7,\allowbreak 15,\allowbreak 5,\allowbreak 13,\allowbreak 3,\allowbreak 10,\allowbreak 1,\allowbreak 8,\allowbreak 11,\allowbreak 14,\allowbreak 6,\allowbreak 12,\allowbreak 4}\par
\noindent \textbf{200}: roots \texttt{3,\allowbreak 4,\allowbreak 6,\allowbreak 5,\allowbreak 6,\allowbreak 7,\allowbreak 6,\allowbreak 7}; order \texttt{1,\allowbreak 2,\allowbreak 7,\allowbreak 4,\allowbreak 5,\allowbreak 8,\allowbreak 3,\allowbreak 6}\par
\noindent \textbf{201}: roots \texttt{3,\allowbreak 4,\allowbreak 3,\allowbreak 4,\allowbreak 6,\allowbreak 5,\allowbreak 4,\allowbreak 6,\allowbreak 7,\allowbreak 6,\allowbreak 7,\allowbreak 6}; order \texttt{7,\allowbreak 3,\allowbreak 12,\allowbreak 1,\allowbreak 4,\allowbreak 10,\allowbreak 2,\allowbreak 6,\allowbreak 8,\allowbreak 11,\allowbreak 5,\allowbreak 9}\par
\noindent \textbf{202}: roots \texttt{4,\allowbreak 6,\allowbreak 4,\allowbreak 6,\allowbreak 4,\allowbreak 6,\allowbreak 7,\allowbreak 6,\allowbreak 7,\allowbreak 6}; order \texttt{5,\allowbreak 10,\allowbreak 3,\allowbreak 8,\allowbreak 1,\allowbreak 6,\allowbreak 9,\allowbreak 4,\allowbreak 7,\allowbreak 2}\par
\noindent \textbf{203}: roots \texttt{4,\allowbreak 6,\allowbreak 3,\allowbreak 4,\allowbreak 6,\allowbreak 4,\allowbreak 6,\allowbreak 7,\allowbreak 6,\allowbreak 7,\allowbreak 6}; order \texttt{3,\allowbreak 6,\allowbreak 11,\allowbreak 4,\allowbreak 9,\allowbreak 1,\allowbreak 7,\allowbreak 10,\allowbreak 5,\allowbreak 8,\allowbreak 2}\par
\noindent \textbf{204}: roots \texttt{4,\allowbreak 6,\allowbreak 3,\allowbreak 4,\allowbreak 6,\allowbreak 5,\allowbreak 6,\allowbreak 7,\allowbreak 6,\allowbreak 7,\allowbreak 6}; order \texttt{3,\allowbreak 11,\allowbreak 4,\allowbreak 9,\allowbreak 1,\allowbreak 6,\allowbreak 7,\allowbreak 10,\allowbreak 5,\allowbreak 8,\allowbreak 2}\par
\noindent \textbf{205}: roots \texttt{4,\allowbreak 2,\allowbreak 3,\allowbreak 4,\allowbreak 6,\allowbreak 5,\allowbreak 4,\allowbreak 5,\allowbreak 6,\allowbreak 7,\allowbreak 6,\allowbreak 7}; order \texttt{7,\allowbreak 2,\allowbreak 3,\allowbreak 4,\allowbreak 11,\allowbreak 8,\allowbreak 1,\allowbreak 6,\allowbreak 9,\allowbreak 12,\allowbreak 5,\allowbreak 10}\par
\noindent \textbf{206}: roots \texttt{2,\allowbreak 3,\allowbreak 4,\allowbreak 6,\allowbreak 4,\allowbreak 6,\allowbreak 5,\allowbreak 4,\allowbreak 5,\allowbreak 6,\allowbreak 7,\allowbreak 6,\allowbreak 7,\allowbreak 6,\allowbreak 7}; order \texttt{1,\allowbreak 8,\allowbreak 2,\allowbreak 5,\allowbreak 14,\allowbreak 3,\allowbreak 9,\allowbreak 12,\allowbreak 7,\allowbreak 10,\allowbreak 15,\allowbreak 6,\allowbreak 13,\allowbreak 4,\allowbreak 11}\par
\noindent \textbf{233}: roots \texttt{2,\allowbreak 6,\allowbreak 4,\allowbreak 5,\allowbreak 6,\allowbreak 5,\allowbreak 6,\allowbreak 5,\allowbreak 6,\allowbreak 7,\allowbreak 6,\allowbreak 7,\allowbreak 6,\allowbreak 7,\allowbreak 6,\allowbreak 7}; order \texttt{1,\allowbreak 15,\allowbreak 8,\allowbreak 13,\allowbreak 3,\allowbreak 6,\allowbreak 11,\allowbreak 4,\allowbreak 9,\allowbreak 16,\allowbreak 7,\allowbreak 14,\allowbreak 5,\allowbreak 12,\allowbreak 2,\allowbreak 10}\par
\noindent \textbf{234}: roots \texttt{6,\allowbreak 5,\allowbreak 6,\allowbreak 5,\allowbreak 6,\allowbreak 5,\allowbreak 7,\allowbreak 6,\allowbreak 7,\allowbreak 6,\allowbreak 7,\allowbreak 6,\allowbreak 7}; order \texttt{12,\allowbreak 6,\allowbreak 10,\allowbreak 4,\allowbreak 8,\allowbreak 2,\allowbreak 13,\allowbreak 5,\allowbreak 11,\allowbreak 3,\allowbreak 9,\allowbreak 1,\allowbreak 7}\par
\noindent \textbf{235}: roots \texttt{2,\allowbreak 6,\allowbreak 3,\allowbreak 5,\allowbreak 6,\allowbreak 5,\allowbreak 6,\allowbreak 5,\allowbreak 7,\allowbreak 6,\allowbreak 7,\allowbreak 6,\allowbreak 7,\allowbreak 6,\allowbreak 7}; order \texttt{1,\allowbreak 14,\allowbreak 3,\allowbreak 8,\allowbreak 12,\allowbreak 6,\allowbreak 10,\allowbreak 4,\allowbreak 15,\allowbreak 7,\allowbreak 13,\allowbreak 5,\allowbreak 11,\allowbreak 2,\allowbreak 9}\par
\noindent \textbf{236}: roots \texttt{2,\allowbreak 6,\allowbreak 5,\allowbreak 6,\allowbreak 4,\allowbreak 5,\allowbreak 6,\allowbreak 5,\allowbreak 7,\allowbreak 6,\allowbreak 7,\allowbreak 6,\allowbreak 7,\allowbreak 6,\allowbreak 7}; order \texttt{1,\allowbreak 14,\allowbreak 5,\allowbreak 8,\allowbreak 12,\allowbreak 6,\allowbreak 10,\allowbreak 3,\allowbreak 15,\allowbreak 7,\allowbreak 13,\allowbreak 4,\allowbreak 11,\allowbreak 2,\allowbreak 9}\par
\noindent \textbf{243}: roots \texttt{4,\allowbreak 5,\allowbreak 4,\allowbreak 6,\allowbreak 5,\allowbreak 7,\allowbreak 6,\allowbreak 7}; order \texttt{3,\allowbreak 7,\allowbreak 5,\allowbreak 1,\allowbreak 2,\allowbreak 8,\allowbreak 4,\allowbreak 6}\par
\noindent \textbf{245}: roots \texttt{4,\allowbreak 5,\allowbreak 6,\allowbreak 5,\allowbreak 7,\allowbreak 6,\allowbreak 7}; order \texttt{6,\allowbreak 4,\allowbreak 1,\allowbreak 2,\allowbreak 7,\allowbreak 3,\allowbreak 5}\par
\noindent \textbf{246}: roots \texttt{6,\allowbreak 3,\allowbreak 6,\allowbreak 5,\allowbreak 4,\allowbreak 5,\allowbreak 6,\allowbreak 4,\allowbreak 6,\allowbreak 5,\allowbreak 7,\allowbreak 6,\allowbreak 7,\allowbreak 6,\allowbreak 7,\allowbreak 6,\allowbreak 7,\allowbreak 6,\allowbreak 7}; order \texttt{8,\allowbreak 2,\allowbreak 18,\allowbreak 5,\allowbreak 16,\allowbreak 10,\allowbreak 14,\allowbreak 6,\allowbreak 12,\allowbreak 4,\allowbreak 19,\allowbreak 9,\allowbreak 17,\allowbreak 7,\allowbreak 15,\allowbreak 3,\allowbreak 13,\allowbreak 1,\allowbreak 11}\par
\noindent \textbf{247}: roots \texttt{6,\allowbreak 2,\allowbreak 3,\allowbreak 2,\allowbreak 6,\allowbreak 3,\allowbreak 6,\allowbreak 5,\allowbreak 4,\allowbreak 5,\allowbreak 4,\allowbreak 5,\allowbreak 6,\allowbreak 4,\allowbreak 6,\allowbreak 5,\allowbreak 7,\allowbreak 6,\allowbreak 7,\allowbreak 6,\allowbreak 7,\allowbreak 6,\allowbreak 7,\allowbreak 6,\allowbreak 7,\allowbreak 6,\allowbreak 7}; order \texttt{14,\allowbreak 4,\allowbreak 6,\allowbreak 2,\allowbreak 11,\allowbreak 3,\allowbreak 26,\allowbreak 9,\allowbreak 24,\allowbreak 16,\allowbreak 22,\allowbreak 12,\allowbreak 20,\allowbreak 10,\allowbreak 18,\allowbreak 8,\allowbreak 27,\allowbreak 15,\allowbreak 25,\allowbreak 13,\allowbreak 23,\allowbreak 7,\allowbreak 21,\allowbreak 5,\allowbreak 19,\allowbreak 1,\allowbreak 17}\par
\noindent \textbf{248}: roots \texttt{6,\allowbreak 2,\allowbreak 6,\allowbreak 5,\allowbreak 4,\allowbreak 5,\allowbreak 4,\allowbreak 6,\allowbreak 4,\allowbreak 6,\allowbreak 5,\allowbreak 7,\allowbreak 6,\allowbreak 7,\allowbreak 6,\allowbreak 7,\allowbreak 6,\allowbreak 7,\allowbreak 6,\allowbreak 7}; order \texttt{9,\allowbreak 2,\allowbreak 7,\allowbreak 19,\allowbreak 5,\allowbreak 17,\allowbreak 11,\allowbreak 15,\allowbreak 6,\allowbreak 13,\allowbreak 4,\allowbreak 20,\allowbreak 10,\allowbreak 18,\allowbreak 8,\allowbreak 16,\allowbreak 3,\allowbreak 14,\allowbreak 1,\allowbreak 12}\par
\noindent \textbf{249}: roots \texttt{3,\allowbreak 2,\allowbreak 6,\allowbreak 5,\allowbreak 4,\allowbreak 6,\allowbreak 4,\allowbreak 6,\allowbreak 5,\allowbreak 7,\allowbreak 6,\allowbreak 7,\allowbreak 6,\allowbreak 7,\allowbreak 6,\allowbreak 7}; order \texttt{7,\allowbreak 2,\allowbreak 5,\allowbreak 15,\allowbreak 1,\allowbreak 13,\allowbreak 9,\allowbreak 11,\allowbreak 4,\allowbreak 16,\allowbreak 8,\allowbreak 14,\allowbreak 6,\allowbreak 12,\allowbreak 3,\allowbreak 10}\par
\noindent \textbf{250}: roots \texttt{2,\allowbreak 6,\allowbreak 5,\allowbreak 4,\allowbreak 6,\allowbreak 5,\allowbreak 7,\allowbreak 6,\allowbreak 7,\allowbreak 6,\allowbreak 7}; order \texttt{1,\allowbreak 4,\allowbreak 10,\allowbreak 6,\allowbreak 8,\allowbreak 3,\allowbreak 11,\allowbreak 5,\allowbreak 9,\allowbreak 2,\allowbreak 7}\par
\noindent \textbf{251}: roots \texttt{2,\allowbreak 3,\allowbreak 6,\allowbreak 5,\allowbreak 6,\allowbreak 5,\allowbreak 7,\allowbreak 6,\allowbreak 7,\allowbreak 6,\allowbreak 7}; order \texttt{1,\allowbreak 2,\allowbreak 10,\allowbreak 6,\allowbreak 8,\allowbreak 4,\allowbreak 11,\allowbreak 5,\allowbreak 9,\allowbreak 3,\allowbreak 7}\par
\noindent \textbf{252}: roots \texttt{2,\allowbreak 6,\allowbreak 4,\allowbreak 6,\allowbreak 5,\allowbreak 6,\allowbreak 5,\allowbreak 6,\allowbreak 5,\allowbreak 7,\allowbreak 6,\allowbreak 7,\allowbreak 6,\allowbreak 7,\allowbreak 6,\allowbreak 7,\allowbreak 6,\allowbreak 7}; order \texttt{1,\allowbreak 17,\allowbreak 9,\allowbreak 15,\allowbreak 3,\allowbreak 13,\allowbreak 7,\allowbreak 11,\allowbreak 5,\allowbreak 18,\allowbreak 8,\allowbreak 16,\allowbreak 6,\allowbreak 14,\allowbreak 4,\allowbreak 12,\allowbreak 2,\allowbreak 10}\par
\noindent \textbf{253}: roots \texttt{4,\allowbreak 3,\allowbreak 6,\allowbreak 4,\allowbreak 6,\allowbreak 4,\allowbreak 5,\allowbreak 6,\allowbreak 5,\allowbreak 6,\allowbreak 5,\allowbreak 7,\allowbreak 6,\allowbreak 7,\allowbreak 6,\allowbreak 7,\allowbreak 6,\allowbreak 7,\allowbreak 6,\allowbreak 7}; order \texttt{19,\allowbreak 2,\allowbreak 17,\allowbreak 6,\allowbreak 11,\allowbreak 4,\allowbreak 9,\allowbreak 15,\allowbreak 1,\allowbreak 13,\allowbreak 7,\allowbreak 20,\allowbreak 10,\allowbreak 18,\allowbreak 8,\allowbreak 16,\allowbreak 5,\allowbreak 14,\allowbreak 3,\allowbreak 12}\par
\noindent \textbf{339}: roots \texttt{3,\allowbreak 6,\allowbreak 4,\allowbreak 3,\allowbreak 6,\allowbreak 8,\allowbreak 7,\allowbreak 8,\allowbreak 7}; order \texttt{4,\allowbreak 9,\allowbreak 1,\allowbreak 7,\allowbreak 3,\allowbreak 5,\allowbreak 8,\allowbreak 2,\allowbreak 6}\par
\noindent \textbf{341}: roots \texttt{3,\allowbreak 7,\allowbreak 6,\allowbreak 4,\allowbreak 6,\allowbreak 7,\allowbreak 8,\allowbreak 7}; order \texttt{4,\allowbreak 8,\allowbreak 1,\allowbreak 5,\allowbreak 6,\allowbreak 3,\allowbreak 7,\allowbreak 2}\par
\noindent \textbf{343}: roots \texttt{5,\allowbreak 4,\allowbreak 6,\allowbreak 4,\allowbreak 6,\allowbreak 5,\allowbreak 6,\allowbreak 8,\allowbreak 7,\allowbreak 8}; order \texttt{9,\allowbreak 4,\allowbreak 6,\allowbreak 2,\allowbreak 7,\allowbreak 1,\allowbreak 5,\allowbreak 10,\allowbreak 3,\allowbreak 8}\par
\noindent \textbf{344}: roots \texttt{7,\allowbreak 6,\allowbreak 5,\allowbreak 6,\allowbreak 7,\allowbreak 8,\allowbreak 7,\allowbreak 6,\allowbreak 7,\allowbreak 8,\allowbreak 7,\allowbreak 8}; order \texttt{11,\allowbreak 8,\allowbreak 3,\allowbreak 4,\allowbreak 9,\allowbreak 12,\allowbreak 7,\allowbreak 2,\allowbreak 5,\allowbreak 10,\allowbreak 1,\allowbreak 6}\par
\noindent \textbf{346}: roots \texttt{5,\allowbreak 7,\allowbreak 4,\allowbreak 6,\allowbreak 5,\allowbreak 6,\allowbreak 8,\allowbreak 7,\allowbreak 6,\allowbreak 7,\allowbreak 8}; order \texttt{3,\allowbreak 9,\allowbreak 5,\allowbreak 6,\allowbreak 1,\allowbreak 10,\allowbreak 4,\allowbreak 8,\allowbreak 11,\allowbreak 2,\allowbreak 7}\par
\noindent \textbf{347}: roots \texttt{3,\allowbreak 4,\allowbreak 6,\allowbreak 4,\allowbreak 6,\allowbreak 7,\allowbreak 8,\allowbreak 7,\allowbreak 8}; order \texttt{8,\allowbreak 4,\allowbreak 1,\allowbreak 2,\allowbreak 5,\allowbreak 9,\allowbreak 3,\allowbreak 6,\allowbreak 7}\par
\noindent \textbf{348}: roots \texttt{6,\allowbreak 5,\allowbreak 6,\allowbreak 5,\allowbreak 6,\allowbreak 7,\allowbreak 6,\allowbreak 7,\allowbreak 8,\allowbreak 7,\allowbreak 8}; order \texttt{10,\allowbreak 4,\allowbreak 7,\allowbreak 2,\allowbreak 5,\allowbreak 11,\allowbreak 3,\allowbreak 8,\allowbreak 1,\allowbreak 6,\allowbreak 9}\par
\noindent \textbf{349}: roots \texttt{7,\allowbreak 4,\allowbreak 6,\allowbreak 5,\allowbreak 6,\allowbreak 8,\allowbreak 6,\allowbreak 7,\allowbreak 6,\allowbreak 8}; order \texttt{9,\allowbreak 2,\allowbreak 7,\allowbreak 4,\allowbreak 5,\allowbreak 8,\allowbreak 3,\allowbreak 10,\allowbreak 1,\allowbreak 6}\par
\noindent \textbf{350}: roots \texttt{7,\allowbreak 4,\allowbreak 3,\allowbreak 4,\allowbreak 6,\allowbreak 8,\allowbreak 6,\allowbreak 8}; order \texttt{3,\allowbreak 4,\allowbreak 7,\allowbreak 2,\allowbreak 5,\allowbreak 8,\allowbreak 1,\allowbreak 6}\par
\noindent \textbf{353}: roots \texttt{3,\allowbreak 2,\allowbreak 6,\allowbreak 3,\allowbreak 5,\allowbreak 4,\allowbreak 6,\allowbreak 7,\allowbreak 6,\allowbreak 8,\allowbreak 7,\allowbreak 8}; order \texttt{2,\allowbreak 4,\allowbreak 6,\allowbreak 1,\allowbreak 11,\allowbreak 9,\allowbreak 5,\allowbreak 7,\allowbreak 12,\allowbreak 3,\allowbreak 8,\allowbreak 10}\par
\noindent \textbf{354}: roots \texttt{7,\allowbreak 4,\allowbreak 3,\allowbreak 4,\allowbreak 6,\allowbreak 8,\allowbreak 7,\allowbreak 6,\allowbreak 8,\allowbreak 7,\allowbreak 8}; order \texttt{10,\allowbreak 3,\allowbreak 4,\allowbreak 8,\allowbreak 2,\allowbreak 5,\allowbreak 11,\allowbreak 7,\allowbreak 9,\allowbreak 1,\allowbreak 6}\par
\noindent \textbf{358}: roots \texttt{3,\allowbreak 2,\allowbreak 7,\allowbreak 4,\allowbreak 5,\allowbreak 6,\allowbreak 8,\allowbreak 6,\allowbreak 7,\allowbreak 8}; order \texttt{2,\allowbreak 4,\allowbreak 1,\allowbreak 8,\allowbreak 5,\allowbreak 6,\allowbreak 9,\allowbreak 10,\allowbreak 3,\allowbreak 7}\par
\noindent \textbf{363}: roots \texttt{7,\allowbreak 6,\allowbreak 4,\allowbreak 6,\allowbreak 5,\allowbreak 6,\allowbreak 7,\allowbreak 8,\allowbreak 7,\allowbreak 8,\allowbreak 7,\allowbreak 8}; order \texttt{11,\allowbreak 3,\allowbreak 5,\allowbreak 6,\allowbreak 12,\allowbreak 9,\allowbreak 4,\allowbreak 7,\allowbreak 2,\allowbreak 10,\allowbreak 1,\allowbreak 8}\par

\end{multicols}


\clearpage
\section{Reduced set of 56 algebraic cubic-order shapes}
\label{app:reduced-shapes}
The critical length $n=173$ is proved directly
by forced-reflection peeling rather than from
the determinant parameter table. For each of
the remaining $56$ shapes, the block lengths
are positive integer determinants by
Theorem~\ref{thm:unimodular-block-factory}.
The excluded original $57$-shape table
remains an optional independent audit.
\begin{multicols}{2}
\scriptsize
\setlength{\parindent}{0pt}
\setlength{\parskip}{2pt}
% 56 residual cubic-order determinant shapes, 784 blocks.
% The excluded seed n=173 has a direct forced-leaf uniqueness proof in FORCED_REFLECTION_PEELING.md.
% Outputs are not given; all block lengths are computed symbolically from Cramer determinants.
\noindent \textbf{82}: roots \texttt{4,\allowbreak 2,\allowbreak 3,\allowbreak 4,\allowbreak 5,\allowbreak 4,\allowbreak 5}; order \texttt{6,\allowbreak 2,\allowbreak 3,\allowbreak 4,\allowbreak 7,\allowbreak 1,\allowbreak 5}\par
\noindent \textbf{116}: roots \texttt{4,\allowbreak 3,\allowbreak 2,\allowbreak 3,\allowbreak 5,\allowbreak 3,\allowbreak 5,\allowbreak 3,\allowbreak 5,\allowbreak 3,\allowbreak 5,\allowbreak 3,\allowbreak 5,\allowbreak 3,\allowbreak 5,\allowbreak 4,\allowbreak 5,\allowbreak 6,\allowbreak 5,\allowbreak 6,\allowbreak 5,\allowbreak 6,\allowbreak 5,\allowbreak 6,\allowbreak 5,\allowbreak 6,\allowbreak 5}; order \texttt{16,\allowbreak 14,\allowbreak 3,\allowbreak 12,\allowbreak 27,\allowbreak 10,\allowbreak 25,\allowbreak 8,\allowbreak 23,\allowbreak 6,\allowbreak 21,\allowbreak 4,\allowbreak 19,\allowbreak 2,\allowbreak 17,\allowbreak 1,\allowbreak 15,\allowbreak 26,\allowbreak 13,\allowbreak 24,\allowbreak 11,\allowbreak 22,\allowbreak 9,\allowbreak 20,\allowbreak 7,\allowbreak 18,\allowbreak 5}\par
\noindent \textbf{118}: roots \texttt{3,\allowbreak 2,\allowbreak 3,\allowbreak 2,\allowbreak 3,\allowbreak 2,\allowbreak 5,\allowbreak 3,\allowbreak 5,\allowbreak 3,\allowbreak 5,\allowbreak 3,\allowbreak 5,\allowbreak 3,\allowbreak 5,\allowbreak 3,\allowbreak 5,\allowbreak 3,\allowbreak 5,\allowbreak 3,\allowbreak 5,\allowbreak 3,\allowbreak 5,\allowbreak 3,\allowbreak 5,\allowbreak 6,\allowbreak 5,\allowbreak 6,\allowbreak 5,\allowbreak 6,\allowbreak 5,\allowbreak 6,\allowbreak 5,\allowbreak 6,\allowbreak 5,\allowbreak 6,\allowbreak 5,\allowbreak 6,\allowbreak 5,\allowbreak 6,\allowbreak 5,\allowbreak 6,\allowbreak 5}; order \texttt{24,\allowbreak 6,\allowbreak 22,\allowbreak 4,\allowbreak 20,\allowbreak 2,\allowbreak 43,\allowbreak 18,\allowbreak 41,\allowbreak 16,\allowbreak 39,\allowbreak 14,\allowbreak 37,\allowbreak 12,\allowbreak 35,\allowbreak 10,\allowbreak 33,\allowbreak 8,\allowbreak 31,\allowbreak 5,\allowbreak 29,\allowbreak 3,\allowbreak 27,\allowbreak 1,\allowbreak 25,\allowbreak 42,\allowbreak 23,\allowbreak 40,\allowbreak 21,\allowbreak 38,\allowbreak 19,\allowbreak 36,\allowbreak 17,\allowbreak 34,\allowbreak 15,\allowbreak 32,\allowbreak 13,\allowbreak 30,\allowbreak 11,\allowbreak 28,\allowbreak 9,\allowbreak 26,\allowbreak 7}\par
\noindent \textbf{120}: roots \texttt{4,\allowbreak 3,\allowbreak 2,\allowbreak 5,\allowbreak 3,\allowbreak 5,\allowbreak 3,\allowbreak 5,\allowbreak 3,\allowbreak 5,\allowbreak 4,\allowbreak 5,\allowbreak 6,\allowbreak 5,\allowbreak 6,\allowbreak 5,\allowbreak 6,\allowbreak 5}; order \texttt{11,\allowbreak 9,\allowbreak 3,\allowbreak 18,\allowbreak 7,\allowbreak 16,\allowbreak 5,\allowbreak 14,\allowbreak 2,\allowbreak 12,\allowbreak 1,\allowbreak 10,\allowbreak 17,\allowbreak 8,\allowbreak 15,\allowbreak 6,\allowbreak 13,\allowbreak 4}\par
\noindent \textbf{121}: roots \texttt{4,\allowbreak 3,\allowbreak 5,\allowbreak 3,\allowbreak 5,\allowbreak 3,\allowbreak 5,\allowbreak 3,\allowbreak 5,\allowbreak 4,\allowbreak 5,\allowbreak 6,\allowbreak 5,\allowbreak 6,\allowbreak 5,\allowbreak 6,\allowbreak 5}; order \texttt{10,\allowbreak 8,\allowbreak 17,\allowbreak 6,\allowbreak 15,\allowbreak 4,\allowbreak 13,\allowbreak 2,\allowbreak 11,\allowbreak 1,\allowbreak 9,\allowbreak 16,\allowbreak 7,\allowbreak 14,\allowbreak 5,\allowbreak 12,\allowbreak 3}\par
\noindent \textbf{123}: roots \texttt{4,\allowbreak 5,\allowbreak 4,\allowbreak 5,\allowbreak 4,\allowbreak 5,\allowbreak 6,\allowbreak 5,\allowbreak 6}; order \texttt{5,\allowbreak 8,\allowbreak 3,\allowbreak 6,\allowbreak 1,\allowbreak 4,\allowbreak 9,\allowbreak 2,\allowbreak 7}\par
\noindent \textbf{125}: roots \texttt{5,\allowbreak 2,\allowbreak 5,\allowbreak 4,\allowbreak 5,\allowbreak 4,\allowbreak 5,\allowbreak 6,\allowbreak 5,\allowbreak 6,\allowbreak 5,\allowbreak 6,\allowbreak 5,\allowbreak 6}; order \texttt{13,\allowbreak 2,\allowbreak 11,\allowbreak 6,\allowbreak 9,\allowbreak 4,\allowbreak 7,\allowbreak 14,\allowbreak 5,\allowbreak 12,\allowbreak 3,\allowbreak 10,\allowbreak 1,\allowbreak 8}\par
\noindent \textbf{126}: roots \texttt{4,\allowbreak 2,\allowbreak 5,\allowbreak 4,\allowbreak 5,\allowbreak 4,\allowbreak 5,\allowbreak 6,\allowbreak 5,\allowbreak 6}; order \texttt{6,\allowbreak 2,\allowbreak 9,\allowbreak 4,\allowbreak 7,\allowbreak 1,\allowbreak 5,\allowbreak 10,\allowbreak 3,\allowbreak 8}\par
\noindent \textbf{127}: roots \texttt{5,\allowbreak 2,\allowbreak 3,\allowbreak 2,\allowbreak 5,\allowbreak 3,\allowbreak 5,\allowbreak 3,\allowbreak 5,\allowbreak 3,\allowbreak 5,\allowbreak 4,\allowbreak 5,\allowbreak 4,\allowbreak 5,\allowbreak 6,\allowbreak 5,\allowbreak 6,\allowbreak 5,\allowbreak 6,\allowbreak 5,\allowbreak 6,\allowbreak 5,\allowbreak 6,\allowbreak 5,\allowbreak 6,\allowbreak 5,\allowbreak 6}; order \texttt{10,\allowbreak 4,\allowbreak 27,\allowbreak 2,\allowbreak 25,\allowbreak 8,\allowbreak 23,\allowbreak 6,\allowbreak 21,\allowbreak 3,\allowbreak 19,\allowbreak 14,\allowbreak 17,\allowbreak 12,\allowbreak 15,\allowbreak 28,\allowbreak 13,\allowbreak 26,\allowbreak 11,\allowbreak 24,\allowbreak 9,\allowbreak 22,\allowbreak 7,\allowbreak 20,\allowbreak 5,\allowbreak 18,\allowbreak 1,\allowbreak 16}\par
\noindent \textbf{129}: roots \texttt{5,\allowbreak 2,\allowbreak 3,\allowbreak 5,\allowbreak 3,\allowbreak 5,\allowbreak 3,\allowbreak 5,\allowbreak 4,\allowbreak 5,\allowbreak 4,\allowbreak 5,\allowbreak 6,\allowbreak 5,\allowbreak 6,\allowbreak 5,\allowbreak 6,\allowbreak 5,\allowbreak 6,\allowbreak 5,\allowbreak 6,\allowbreak 5,\allowbreak 6}; order \texttt{22,\allowbreak 2,\allowbreak 7,\allowbreak 20,\allowbreak 5,\allowbreak 18,\allowbreak 3,\allowbreak 16,\allowbreak 11,\allowbreak 14,\allowbreak 9,\allowbreak 12,\allowbreak 23,\allowbreak 10,\allowbreak 21,\allowbreak 8,\allowbreak 19,\allowbreak 6,\allowbreak 17,\allowbreak 4,\allowbreak 15,\allowbreak 1,\allowbreak 13}\par
\noindent \textbf{176}: roots \texttt{2,\allowbreak 3,\allowbreak 5,\allowbreak 4,\allowbreak 3,\allowbreak 5,\allowbreak 4,\allowbreak 6,\allowbreak 5,\allowbreak 6,\allowbreak 5,\allowbreak 6,\allowbreak 7,\allowbreak 6}; order \texttt{1,\allowbreak 5,\allowbreak 11,\allowbreak 2,\allowbreak 9,\allowbreak 7,\allowbreak 14,\allowbreak 4,\allowbreak 12,\allowbreak 6,\allowbreak 10,\allowbreak 3,\allowbreak 8,\allowbreak 13}\par
\noindent \textbf{177}: roots \texttt{5,\allowbreak 2,\allowbreak 3,\allowbreak 5,\allowbreak 3,\allowbreak 5,\allowbreak 3,\allowbreak 5,\allowbreak 4,\allowbreak 5,\allowbreak 4,\allowbreak 6,\allowbreak 5,\allowbreak 6,\allowbreak 5,\allowbreak 6,\allowbreak 5,\allowbreak 6,\allowbreak 5,\allowbreak 6,\allowbreak 5,\allowbreak 6}; order \texttt{21,\allowbreak 2,\allowbreak 7,\allowbreak 19,\allowbreak 5,\allowbreak 17,\allowbreak 3,\allowbreak 15,\allowbreak 11,\allowbreak 13,\allowbreak 9,\allowbreak 22,\allowbreak 10,\allowbreak 20,\allowbreak 8,\allowbreak 18,\allowbreak 6,\allowbreak 16,\allowbreak 4,\allowbreak 14,\allowbreak 1,\allowbreak 12}\par
\noindent \textbf{179}: roots \texttt{5,\allowbreak 3,\allowbreak 4,\allowbreak 3,\allowbreak 4,\allowbreak 5,\allowbreak 4,\allowbreak 6,\allowbreak 4,\allowbreak 6,\allowbreak 5,\allowbreak 6,\allowbreak 5,\allowbreak 6,\allowbreak 7,\allowbreak 6,\allowbreak 7}; order \texttt{13,\allowbreak 4,\allowbreak 9,\allowbreak 2,\allowbreak 7,\allowbreak 11,\allowbreak 5,\allowbreak 16,\allowbreak 3,\allowbreak 14,\allowbreak 6,\allowbreak 12,\allowbreak 1,\allowbreak 10,\allowbreak 17,\allowbreak 8,\allowbreak 15}\par
\noindent \textbf{180}: roots \texttt{3,\allowbreak 5,\allowbreak 3,\allowbreak 5,\allowbreak 3,\allowbreak 5,\allowbreak 3,\allowbreak 5,\allowbreak 4,\allowbreak 6,\allowbreak 5,\allowbreak 6,\allowbreak 5,\allowbreak 6,\allowbreak 5,\allowbreak 6,\allowbreak 5,\allowbreak 6}; order \texttt{7,\allowbreak 17,\allowbreak 5,\allowbreak 15,\allowbreak 3,\allowbreak 13,\allowbreak 1,\allowbreak 11,\allowbreak 9,\allowbreak 18,\allowbreak 8,\allowbreak 16,\allowbreak 6,\allowbreak 14,\allowbreak 4,\allowbreak 12,\allowbreak 2,\allowbreak 10}\par
\noindent \textbf{181}: roots \texttt{5,\allowbreak 2,\allowbreak 5,\allowbreak 4,\allowbreak 5,\allowbreak 4,\allowbreak 6,\allowbreak 5,\allowbreak 6,\allowbreak 5,\allowbreak 6,\allowbreak 5,\allowbreak 6}; order \texttt{12,\allowbreak 2,\allowbreak 10,\allowbreak 6,\allowbreak 8,\allowbreak 4,\allowbreak 13,\allowbreak 5,\allowbreak 11,\allowbreak 3,\allowbreak 9,\allowbreak 1,\allowbreak 7}\par
\noindent \textbf{184}: roots \texttt{2,\allowbreak 4,\allowbreak 6,\allowbreak 4,\allowbreak 6,\allowbreak 7,\allowbreak 6}; order \texttt{1,\allowbreak 4,\allowbreak 7,\allowbreak 2,\allowbreak 5,\allowbreak 6,\allowbreak 3}\par
\noindent \textbf{185}: roots \texttt{5,\allowbreak 3,\allowbreak 4,\allowbreak 6,\allowbreak 4,\allowbreak 6,\allowbreak 5,\allowbreak 6,\allowbreak 7}; order \texttt{2,\allowbreak 5,\allowbreak 7,\allowbreak 3,\allowbreak 8,\allowbreak 1,\allowbreak 6,\allowbreak 9,\allowbreak 4}\par
\noindent \textbf{188}: roots \texttt{5,\allowbreak 4,\allowbreak 2,\allowbreak 4,\allowbreak 6,\allowbreak 4,\allowbreak 6,\allowbreak 5,\allowbreak 6,\allowbreak 7,\allowbreak 6,\allowbreak 7}; order \texttt{3,\allowbreak 6,\allowbreak 8,\allowbreak 4,\allowbreak 11,\allowbreak 2,\allowbreak 9,\allowbreak 1,\allowbreak 7,\allowbreak 12,\allowbreak 5,\allowbreak 10}\par
\noindent \textbf{194}: roots \texttt{2,\allowbreak 3,\allowbreak 5,\allowbreak 3,\allowbreak 5,\allowbreak 6,\allowbreak 5,\allowbreak 6,\allowbreak 5,\allowbreak 6}; order \texttt{1,\allowbreak 4,\allowbreak 9,\allowbreak 2,\allowbreak 7,\allowbreak 10,\allowbreak 5,\allowbreak 8,\allowbreak 3,\allowbreak 6}\par
\noindent \textbf{197}: roots \texttt{5,\allowbreak 3,\allowbreak 6,\allowbreak 4,\allowbreak 6,\allowbreak 5,\allowbreak 6,\allowbreak 7,\allowbreak 6}; order \texttt{4,\allowbreak 2,\allowbreak 9,\allowbreak 6,\allowbreak 7,\allowbreak 1,\allowbreak 5,\allowbreak 8,\allowbreak 3}\par
\noindent \textbf{198}: roots \texttt{4,\allowbreak 2,\allowbreak 5,\allowbreak 4,\allowbreak 6,\allowbreak 5,\allowbreak 4,\allowbreak 5,\allowbreak 6,\allowbreak 5,\allowbreak 6,\allowbreak 7,\allowbreak 6,\allowbreak 7}; order \texttt{7,\allowbreak 2,\allowbreak 10,\allowbreak 4,\allowbreak 13,\allowbreak 8,\allowbreak 1,\allowbreak 6,\allowbreak 11,\allowbreak 3,\allowbreak 9,\allowbreak 14,\allowbreak 5,\allowbreak 12}\par
\noindent \textbf{199}: roots \texttt{4,\allowbreak 2,\allowbreak 4,\allowbreak 6,\allowbreak 4,\allowbreak 6,\allowbreak 4,\allowbreak 5,\allowbreak 4,\allowbreak 5,\allowbreak 6,\allowbreak 7,\allowbreak 6,\allowbreak 7,\allowbreak 6}; order \texttt{9,\allowbreak 2,\allowbreak 7,\allowbreak 15,\allowbreak 5,\allowbreak 13,\allowbreak 3,\allowbreak 10,\allowbreak 1,\allowbreak 8,\allowbreak 11,\allowbreak 14,\allowbreak 6,\allowbreak 12,\allowbreak 4}\par
\noindent \textbf{200}: roots \texttt{3,\allowbreak 4,\allowbreak 6,\allowbreak 5,\allowbreak 6,\allowbreak 7,\allowbreak 6,\allowbreak 7}; order \texttt{1,\allowbreak 2,\allowbreak 7,\allowbreak 4,\allowbreak 5,\allowbreak 8,\allowbreak 3,\allowbreak 6}\par
\noindent \textbf{201}: roots \texttt{3,\allowbreak 4,\allowbreak 3,\allowbreak 4,\allowbreak 6,\allowbreak 5,\allowbreak 4,\allowbreak 6,\allowbreak 7,\allowbreak 6,\allowbreak 7,\allowbreak 6}; order \texttt{7,\allowbreak 3,\allowbreak 12,\allowbreak 1,\allowbreak 4,\allowbreak 10,\allowbreak 2,\allowbreak 6,\allowbreak 8,\allowbreak 11,\allowbreak 5,\allowbreak 9}\par
\noindent \textbf{202}: roots \texttt{4,\allowbreak 6,\allowbreak 4,\allowbreak 6,\allowbreak 4,\allowbreak 6,\allowbreak 7,\allowbreak 6,\allowbreak 7,\allowbreak 6}; order \texttt{5,\allowbreak 10,\allowbreak 3,\allowbreak 8,\allowbreak 1,\allowbreak 6,\allowbreak 9,\allowbreak 4,\allowbreak 7,\allowbreak 2}\par
\noindent \textbf{203}: roots \texttt{4,\allowbreak 6,\allowbreak 3,\allowbreak 4,\allowbreak 6,\allowbreak 4,\allowbreak 6,\allowbreak 7,\allowbreak 6,\allowbreak 7,\allowbreak 6}; order \texttt{3,\allowbreak 6,\allowbreak 11,\allowbreak 4,\allowbreak 9,\allowbreak 1,\allowbreak 7,\allowbreak 10,\allowbreak 5,\allowbreak 8,\allowbreak 2}\par
\noindent \textbf{204}: roots \texttt{4,\allowbreak 6,\allowbreak 3,\allowbreak 4,\allowbreak 6,\allowbreak 5,\allowbreak 6,\allowbreak 7,\allowbreak 6,\allowbreak 7,\allowbreak 6}; order \texttt{3,\allowbreak 11,\allowbreak 4,\allowbreak 9,\allowbreak 1,\allowbreak 6,\allowbreak 7,\allowbreak 10,\allowbreak 5,\allowbreak 8,\allowbreak 2}\par
\noindent \textbf{205}: roots \texttt{4,\allowbreak 2,\allowbreak 3,\allowbreak 4,\allowbreak 6,\allowbreak 5,\allowbreak 4,\allowbreak 5,\allowbreak 6,\allowbreak 7,\allowbreak 6,\allowbreak 7}; order \texttt{7,\allowbreak 2,\allowbreak 3,\allowbreak 4,\allowbreak 11,\allowbreak 8,\allowbreak 1,\allowbreak 6,\allowbreak 9,\allowbreak 12,\allowbreak 5,\allowbreak 10}\par
\noindent \textbf{206}: roots \texttt{2,\allowbreak 3,\allowbreak 4,\allowbreak 6,\allowbreak 4,\allowbreak 6,\allowbreak 5,\allowbreak 4,\allowbreak 5,\allowbreak 6,\allowbreak 7,\allowbreak 6,\allowbreak 7,\allowbreak 6,\allowbreak 7}; order \texttt{1,\allowbreak 8,\allowbreak 2,\allowbreak 5,\allowbreak 14,\allowbreak 3,\allowbreak 9,\allowbreak 12,\allowbreak 7,\allowbreak 10,\allowbreak 15,\allowbreak 6,\allowbreak 13,\allowbreak 4,\allowbreak 11}\par
\noindent \textbf{233}: roots \texttt{2,\allowbreak 6,\allowbreak 4,\allowbreak 5,\allowbreak 6,\allowbreak 5,\allowbreak 6,\allowbreak 5,\allowbreak 6,\allowbreak 7,\allowbreak 6,\allowbreak 7,\allowbreak 6,\allowbreak 7,\allowbreak 6,\allowbreak 7}; order \texttt{1,\allowbreak 15,\allowbreak 8,\allowbreak 13,\allowbreak 3,\allowbreak 6,\allowbreak 11,\allowbreak 4,\allowbreak 9,\allowbreak 16,\allowbreak 7,\allowbreak 14,\allowbreak 5,\allowbreak 12,\allowbreak 2,\allowbreak 10}\par
\noindent \textbf{234}: roots \texttt{6,\allowbreak 5,\allowbreak 6,\allowbreak 5,\allowbreak 6,\allowbreak 5,\allowbreak 7,\allowbreak 6,\allowbreak 7,\allowbreak 6,\allowbreak 7,\allowbreak 6,\allowbreak 7}; order \texttt{12,\allowbreak 6,\allowbreak 10,\allowbreak 4,\allowbreak 8,\allowbreak 2,\allowbreak 13,\allowbreak 5,\allowbreak 11,\allowbreak 3,\allowbreak 9,\allowbreak 1,\allowbreak 7}\par
\noindent \textbf{235}: roots \texttt{2,\allowbreak 6,\allowbreak 3,\allowbreak 5,\allowbreak 6,\allowbreak 5,\allowbreak 6,\allowbreak 5,\allowbreak 7,\allowbreak 6,\allowbreak 7,\allowbreak 6,\allowbreak 7,\allowbreak 6,\allowbreak 7}; order \texttt{1,\allowbreak 14,\allowbreak 3,\allowbreak 8,\allowbreak 12,\allowbreak 6,\allowbreak 10,\allowbreak 4,\allowbreak 15,\allowbreak 7,\allowbreak 13,\allowbreak 5,\allowbreak 11,\allowbreak 2,\allowbreak 9}\par
\noindent \textbf{236}: roots \texttt{2,\allowbreak 6,\allowbreak 5,\allowbreak 6,\allowbreak 4,\allowbreak 5,\allowbreak 6,\allowbreak 5,\allowbreak 7,\allowbreak 6,\allowbreak 7,\allowbreak 6,\allowbreak 7,\allowbreak 6,\allowbreak 7}; order \texttt{1,\allowbreak 14,\allowbreak 5,\allowbreak 8,\allowbreak 12,\allowbreak 6,\allowbreak 10,\allowbreak 3,\allowbreak 15,\allowbreak 7,\allowbreak 13,\allowbreak 4,\allowbreak 11,\allowbreak 2,\allowbreak 9}\par
\noindent \textbf{243}: roots \texttt{4,\allowbreak 5,\allowbreak 4,\allowbreak 6,\allowbreak 5,\allowbreak 7,\allowbreak 6,\allowbreak 7}; order \texttt{3,\allowbreak 7,\allowbreak 5,\allowbreak 1,\allowbreak 2,\allowbreak 8,\allowbreak 4,\allowbreak 6}\par
\noindent \textbf{245}: roots \texttt{4,\allowbreak 5,\allowbreak 6,\allowbreak 5,\allowbreak 7,\allowbreak 6,\allowbreak 7}; order \texttt{6,\allowbreak 4,\allowbreak 1,\allowbreak 2,\allowbreak 7,\allowbreak 3,\allowbreak 5}\par
\noindent \textbf{246}: roots \texttt{6,\allowbreak 3,\allowbreak 6,\allowbreak 5,\allowbreak 4,\allowbreak 5,\allowbreak 6,\allowbreak 4,\allowbreak 6,\allowbreak 5,\allowbreak 7,\allowbreak 6,\allowbreak 7,\allowbreak 6,\allowbreak 7,\allowbreak 6,\allowbreak 7,\allowbreak 6,\allowbreak 7}; order \texttt{8,\allowbreak 2,\allowbreak 18,\allowbreak 5,\allowbreak 16,\allowbreak 10,\allowbreak 14,\allowbreak 6,\allowbreak 12,\allowbreak 4,\allowbreak 19,\allowbreak 9,\allowbreak 17,\allowbreak 7,\allowbreak 15,\allowbreak 3,\allowbreak 13,\allowbreak 1,\allowbreak 11}\par
\noindent \textbf{247}: roots \texttt{6,\allowbreak 2,\allowbreak 3,\allowbreak 2,\allowbreak 6,\allowbreak 3,\allowbreak 6,\allowbreak 5,\allowbreak 4,\allowbreak 5,\allowbreak 4,\allowbreak 5,\allowbreak 6,\allowbreak 4,\allowbreak 6,\allowbreak 5,\allowbreak 7,\allowbreak 6,\allowbreak 7,\allowbreak 6,\allowbreak 7,\allowbreak 6,\allowbreak 7,\allowbreak 6,\allowbreak 7,\allowbreak 6,\allowbreak 7}; order \texttt{14,\allowbreak 4,\allowbreak 6,\allowbreak 2,\allowbreak 11,\allowbreak 3,\allowbreak 26,\allowbreak 9,\allowbreak 24,\allowbreak 16,\allowbreak 22,\allowbreak 12,\allowbreak 20,\allowbreak 10,\allowbreak 18,\allowbreak 8,\allowbreak 27,\allowbreak 15,\allowbreak 25,\allowbreak 13,\allowbreak 23,\allowbreak 7,\allowbreak 21,\allowbreak 5,\allowbreak 19,\allowbreak 1,\allowbreak 17}\par
\noindent \textbf{248}: roots \texttt{6,\allowbreak 2,\allowbreak 6,\allowbreak 5,\allowbreak 4,\allowbreak 5,\allowbreak 4,\allowbreak 6,\allowbreak 4,\allowbreak 6,\allowbreak 5,\allowbreak 7,\allowbreak 6,\allowbreak 7,\allowbreak 6,\allowbreak 7,\allowbreak 6,\allowbreak 7,\allowbreak 6,\allowbreak 7}; order \texttt{9,\allowbreak 2,\allowbreak 7,\allowbreak 19,\allowbreak 5,\allowbreak 17,\allowbreak 11,\allowbreak 15,\allowbreak 6,\allowbreak 13,\allowbreak 4,\allowbreak 20,\allowbreak 10,\allowbreak 18,\allowbreak 8,\allowbreak 16,\allowbreak 3,\allowbreak 14,\allowbreak 1,\allowbreak 12}\par
\noindent \textbf{249}: roots \texttt{3,\allowbreak 2,\allowbreak 6,\allowbreak 5,\allowbreak 4,\allowbreak 6,\allowbreak 4,\allowbreak 6,\allowbreak 5,\allowbreak 7,\allowbreak 6,\allowbreak 7,\allowbreak 6,\allowbreak 7,\allowbreak 6,\allowbreak 7}; order \texttt{7,\allowbreak 2,\allowbreak 5,\allowbreak 15,\allowbreak 1,\allowbreak 13,\allowbreak 9,\allowbreak 11,\allowbreak 4,\allowbreak 16,\allowbreak 8,\allowbreak 14,\allowbreak 6,\allowbreak 12,\allowbreak 3,\allowbreak 10}\par
\noindent \textbf{250}: roots \texttt{2,\allowbreak 6,\allowbreak 5,\allowbreak 4,\allowbreak 6,\allowbreak 5,\allowbreak 7,\allowbreak 6,\allowbreak 7,\allowbreak 6,\allowbreak 7}; order \texttt{1,\allowbreak 4,\allowbreak 10,\allowbreak 6,\allowbreak 8,\allowbreak 3,\allowbreak 11,\allowbreak 5,\allowbreak 9,\allowbreak 2,\allowbreak 7}\par
\noindent \textbf{251}: roots \texttt{2,\allowbreak 3,\allowbreak 6,\allowbreak 5,\allowbreak 6,\allowbreak 5,\allowbreak 7,\allowbreak 6,\allowbreak 7,\allowbreak 6,\allowbreak 7}; order \texttt{1,\allowbreak 2,\allowbreak 10,\allowbreak 6,\allowbreak 8,\allowbreak 4,\allowbreak 11,\allowbreak 5,\allowbreak 9,\allowbreak 3,\allowbreak 7}\par
\noindent \textbf{252}: roots \texttt{2,\allowbreak 6,\allowbreak 4,\allowbreak 6,\allowbreak 5,\allowbreak 6,\allowbreak 5,\allowbreak 6,\allowbreak 5,\allowbreak 7,\allowbreak 6,\allowbreak 7,\allowbreak 6,\allowbreak 7,\allowbreak 6,\allowbreak 7,\allowbreak 6,\allowbreak 7}; order \texttt{1,\allowbreak 17,\allowbreak 9,\allowbreak 15,\allowbreak 3,\allowbreak 13,\allowbreak 7,\allowbreak 11,\allowbreak 5,\allowbreak 18,\allowbreak 8,\allowbreak 16,\allowbreak 6,\allowbreak 14,\allowbreak 4,\allowbreak 12,\allowbreak 2,\allowbreak 10}\par
\noindent \textbf{253}: roots \texttt{4,\allowbreak 3,\allowbreak 6,\allowbreak 4,\allowbreak 6,\allowbreak 4,\allowbreak 5,\allowbreak 6,\allowbreak 5,\allowbreak 6,\allowbreak 5,\allowbreak 7,\allowbreak 6,\allowbreak 7,\allowbreak 6,\allowbreak 7,\allowbreak 6,\allowbreak 7,\allowbreak 6,\allowbreak 7}; order \texttt{19,\allowbreak 2,\allowbreak 17,\allowbreak 6,\allowbreak 11,\allowbreak 4,\allowbreak 9,\allowbreak 15,\allowbreak 1,\allowbreak 13,\allowbreak 7,\allowbreak 20,\allowbreak 10,\allowbreak 18,\allowbreak 8,\allowbreak 16,\allowbreak 5,\allowbreak 14,\allowbreak 3,\allowbreak 12}\par
\noindent \textbf{339}: roots \texttt{3,\allowbreak 6,\allowbreak 4,\allowbreak 3,\allowbreak 6,\allowbreak 8,\allowbreak 7,\allowbreak 8,\allowbreak 7}; order \texttt{4,\allowbreak 9,\allowbreak 1,\allowbreak 7,\allowbreak 3,\allowbreak 5,\allowbreak 8,\allowbreak 2,\allowbreak 6}\par
\noindent \textbf{341}: roots \texttt{3,\allowbreak 7,\allowbreak 6,\allowbreak 4,\allowbreak 6,\allowbreak 7,\allowbreak 8,\allowbreak 7}; order \texttt{4,\allowbreak 8,\allowbreak 1,\allowbreak 5,\allowbreak 6,\allowbreak 3,\allowbreak 7,\allowbreak 2}\par
\noindent \textbf{343}: roots \texttt{5,\allowbreak 4,\allowbreak 6,\allowbreak 4,\allowbreak 6,\allowbreak 5,\allowbreak 6,\allowbreak 8,\allowbreak 7,\allowbreak 8}; order \texttt{9,\allowbreak 4,\allowbreak 6,\allowbreak 2,\allowbreak 7,\allowbreak 1,\allowbreak 5,\allowbreak 10,\allowbreak 3,\allowbreak 8}\par
\noindent \textbf{344}: roots \texttt{7,\allowbreak 6,\allowbreak 5,\allowbreak 6,\allowbreak 7,\allowbreak 8,\allowbreak 7,\allowbreak 6,\allowbreak 7,\allowbreak 8,\allowbreak 7,\allowbreak 8}; order \texttt{11,\allowbreak 8,\allowbreak 3,\allowbreak 4,\allowbreak 9,\allowbreak 12,\allowbreak 7,\allowbreak 2,\allowbreak 5,\allowbreak 10,\allowbreak 1,\allowbreak 6}\par
\noindent \textbf{346}: roots \texttt{5,\allowbreak 7,\allowbreak 4,\allowbreak 6,\allowbreak 5,\allowbreak 6,\allowbreak 8,\allowbreak 7,\allowbreak 6,\allowbreak 7,\allowbreak 8}; order \texttt{3,\allowbreak 9,\allowbreak 5,\allowbreak 6,\allowbreak 1,\allowbreak 10,\allowbreak 4,\allowbreak 8,\allowbreak 11,\allowbreak 2,\allowbreak 7}\par
\noindent \textbf{347}: roots \texttt{3,\allowbreak 4,\allowbreak 6,\allowbreak 4,\allowbreak 6,\allowbreak 7,\allowbreak 8,\allowbreak 7,\allowbreak 8}; order \texttt{8,\allowbreak 4,\allowbreak 1,\allowbreak 2,\allowbreak 5,\allowbreak 9,\allowbreak 3,\allowbreak 6,\allowbreak 7}\par
\noindent \textbf{348}: roots \texttt{6,\allowbreak 5,\allowbreak 6,\allowbreak 5,\allowbreak 6,\allowbreak 7,\allowbreak 6,\allowbreak 7,\allowbreak 8,\allowbreak 7,\allowbreak 8}; order \texttt{10,\allowbreak 4,\allowbreak 7,\allowbreak 2,\allowbreak 5,\allowbreak 11,\allowbreak 3,\allowbreak 8,\allowbreak 1,\allowbreak 6,\allowbreak 9}\par
\noindent \textbf{349}: roots \texttt{7,\allowbreak 4,\allowbreak 6,\allowbreak 5,\allowbreak 6,\allowbreak 8,\allowbreak 6,\allowbreak 7,\allowbreak 6,\allowbreak 8}; order \texttt{9,\allowbreak 2,\allowbreak 7,\allowbreak 4,\allowbreak 5,\allowbreak 8,\allowbreak 3,\allowbreak 10,\allowbreak 1,\allowbreak 6}\par
\noindent \textbf{350}: roots \texttt{7,\allowbreak 4,\allowbreak 3,\allowbreak 4,\allowbreak 6,\allowbreak 8,\allowbreak 6,\allowbreak 8}; order \texttt{3,\allowbreak 4,\allowbreak 7,\allowbreak 2,\allowbreak 5,\allowbreak 8,\allowbreak 1,\allowbreak 6}\par
\noindent \textbf{353}: roots \texttt{3,\allowbreak 2,\allowbreak 6,\allowbreak 3,\allowbreak 5,\allowbreak 4,\allowbreak 6,\allowbreak 7,\allowbreak 6,\allowbreak 8,\allowbreak 7,\allowbreak 8}; order \texttt{2,\allowbreak 4,\allowbreak 6,\allowbreak 1,\allowbreak 11,\allowbreak 9,\allowbreak 5,\allowbreak 7,\allowbreak 12,\allowbreak 3,\allowbreak 8,\allowbreak 10}\par
\noindent \textbf{354}: roots \texttt{7,\allowbreak 4,\allowbreak 3,\allowbreak 4,\allowbreak 6,\allowbreak 8,\allowbreak 7,\allowbreak 6,\allowbreak 8,\allowbreak 7,\allowbreak 8}; order \texttt{10,\allowbreak 3,\allowbreak 4,\allowbreak 8,\allowbreak 2,\allowbreak 5,\allowbreak 11,\allowbreak 7,\allowbreak 9,\allowbreak 1,\allowbreak 6}\par
\noindent \textbf{358}: roots \texttt{3,\allowbreak 2,\allowbreak 7,\allowbreak 4,\allowbreak 5,\allowbreak 6,\allowbreak 8,\allowbreak 6,\allowbreak 7,\allowbreak 8}; order \texttt{2,\allowbreak 4,\allowbreak 1,\allowbreak 8,\allowbreak 5,\allowbreak 6,\allowbreak 9,\allowbreak 10,\allowbreak 3,\allowbreak 7}\par
\noindent \textbf{363}: roots \texttt{7,\allowbreak 6,\allowbreak 4,\allowbreak 6,\allowbreak 5,\allowbreak 6,\allowbreak 7,\allowbreak 8,\allowbreak 7,\allowbreak 8,\allowbreak 7,\allowbreak 8}; order \texttt{11,\allowbreak 3,\allowbreak 5,\allowbreak 6,\allowbreak 12,\allowbreak 9,\allowbreak 4,\allowbreak 7,\allowbreak 2,\allowbreak 10,\allowbreak 1,\allowbreak 8}\par

\end{multicols}

\end{document}
